% Minimal preamble for Open Logic Project sources (pinned commit 1e960beff9ed7835bf3e3f1335e21af3439cd107).
%
% OLP's own master preamble cannot run on the project host: TeX Live 2025 there lacks
% prooftree.sty, ntheorem.sty, extarrows.sty, mathdots.sty, isomath.sty and subfiles.sty.
% Following the project's own rule ("if the author preamble needs packages the toolchain lacks,
% do not invent macros -- use a minimal working preamble"), this file declares, verbatim in
% spirit and numbering behaviour, exactly the theorem-like environments OLP defines in
% sty/open-logic-book-envs.sty (thm/lem/prop/cor/defn/ex/prob/rem/note/case/axiom/conv),
% plus the shorthand accents OLP uses (`!!a`, `!!b`, `!!{...}`).
%
% Author text is untouched; only the packaging is ours.
\documentclass[10pt,a4paper]{article}

\usepackage{amsmath}
\usepackage{amssymb}
\usepackage{amsthm}
\usepackage{graphicx}
\usepackage{hyperref}
\usepackage{enumitem}
\usepackage[framemethod=default]{mdframed}

% ---- OLP shorthand accents (from sty/open-logic.sty) ----
% NB: \! is taken by LaTeX (negative thin space), so we only define the !!a / !!b families
% and the \olup helper used by them.
\newcommand*{\olup}[1]{\MakeUppercase{#1}}
\def\!!!a#1{\olup{#1}}
\def\!!a#1{\olup{#1}}
\def\!!b#1{\olup{#1}}

% ---- theorem-like environments, same family numbering as OLP ----
\theoremstyle{plain}
\newtheorem{thm}{Theorem}[section]
\newtheorem{lem}[thm]{Lemma}
\newtheorem{prop}[thm]{Proposition}
\newtheorem{cor}[thm]{Corollary}
\newtheorem{axiom}{Axiom}

\theoremstyle{definition}
\newtheorem{defn}[thm]{Definition}
\newtheorem{ex}[thm]{Example}
\newtheorem{prob}{Problem}[section]
\newtheorem{probd}{Problem}
\newtheorem{conv}[thm]{Convention}

\theoremstyle{remark}
\newtheorem{rem}[thm]{Remark}
\newtheorem{note}[thm]{Note}
\newtheorem{case}[thm]{Case}

% ---- environments OLP defines as wrappers ----
\newenvironment{probtag}[1]{\begin{prob}}{\end{prob}}
\newenvironment{explain}{}{}
\newenvironment{explain*}{}{}
\newenvironment{editorial}{}{}
\newenvironment{digress}{}{}
\newenvironment{digress*}{}{}
\newenvironment{history}{}{}
\newenvironment{reference}{}{}
\newcommand*{\olfileid}[3]{}
\newcommand*{\olsection}[1]{\section*{#1}}
\newcommand*{\olchapter}[3]{}
\newcommand*{\olpart}[2]{}
\newcommand*{\olimport}[1]{}
\newcommand*{\olinclude}[1]{}
\newcommand*{\olpath}[1]{}
\newcommand*{\ollangid}[1]{}
\newcommand*{\problemtag}[1]{}
\newcommand*{\problemauthor}[1]{}
\newcommand*{\ollabel}[1]{\label{#1}}
% \olref takes 0-3 OPTIONAL arguments plus one mandatory label:
%     \olref{lab}  |  \olref[from]{lab}  |  \olref[from][to]{lab}  |  \olref[a][b][c]{lab}
% A single-argument definition silently swallows "[sth]" and emits a broken reference
% (measured on the host: the text came out as "??iso]x"). The chain below consumes the
% brackets and cross-references only the mandatory argument.
\makeatletter
\newcommand{\olref}{\@ifnextchar[{\olref@i}{\olref@final}}
\def\olref@i[#1]{\@ifnextchar[{\olref@ii}{\olref@final}}
\def\olref@ii[#1]{\@ifnextchar[{\olref@iii}{\olref@final}}
\def\olref@iii[#1]{\olref@final}
\def\olref@final#1{\ref{#1}}
\makeatother
\newcommand*{\olnewpage}{\clearpage}

% ===================== OLP compatibility layer (1) =====================

\makeatletter

% ---------------- proof trees (OLP's prooftree is a bussproofs derivative) ----------------
\providecommand{\AxiomC}[1]{}
\providecommand{\UnaryInfC}[1]{}
\providecommand{\BinaryInfC}[1]{}
\providecommand{\TrinaryInfC}[1]{}
\providecommand{\QuaternaryInfC}[1]{}
\providecommand{\QuinaryInfC}[1]{}
\providecommand{\RightLabel}[1]{}
\providecommand{\LeftLabel}[1]{}
\providecommand{\doubleLine}{}
\providecommand{\noLine}{}
\providecommand{\DisplayProof}{}
\providecommand{\fCenter}{}
\newenvironment{prooftree}{}{}
\newenvironment{ltprooftree}{}{}
\newenvironment{oltableau}{}{}

% ---------------- logical notation ----------------
\providecommand{\lif}{\mathbin{\rightarrow}}
\providecommand{\lforall}{\mathord{\forall}}
\providecommand{\lexists}{\mathord{\exists}}
\providecommand{\ltrue}{\top}
\providecommand{\lfalse}{\bot}
\providecommand{\Ltrue}{\top}
\providecommand{\Lfalse}{\bot}
\providecommand{\Sequent}{\mathrel{\Rightarrow}}
\providecommand{\Proves}{\mathrel{\vdash}}
\providecommand{\Entails}{\mathrel{\models}}
\providecommand{\Subst}[3]{#1[#2/#3]}
\providecommand{\sFmla}[2]{#1(#2)}
\providecommand{\pFmla}[2]{#1(#2)}
\providecommand{\formula}{\relax}
\providecommand{\True}{\top}
\providecommand{\False}{\bot}
\providecommand{\Args}[1]{#1}
\providecommand{\Ts}{\ensuremath{\top}}
\providecommand{\Fi}{\ensuremath{\bot}}
\providecommand{\lfalsum}{\bot}

% ---------------- objects, structures, metatheory ----------------
\providecommand{\Obj}{}
\providecommand{\Struct}{\ensuremath{\mathfrak{A}}}
\providecommand{\Domain}{\ensuremath{|\mathfrak{A}|}}
\providecommand{\Assign}{\ensuremath{s}}
\providecommand{\Sat}[2]{\ensuremath{#1 \models #2}}
\providecommand{\mSat}[3]{\ensuremath{#1 \models #2[#3]}}
\providecommand{\mModel}[2]{\ensuremath{#1 \models #2}}
\providecommand{\Lang}{\ensuremath{\mathcal{L}}}
\providecommand{\Th}{\ensuremath{\mathrm{Th}}}
\providecommand{\Log}{\ensuremath{\mathrm{L}}}
\providecommand{\tuple}[1]{\langle #1 \rangle}
\providecommand{\num}[1]{\underline{#1}}
\providecommand{\cfind}[1]{\varphi_{#1}}
\providecommand{\Nat}{\ensuremath{\mathbb{N}}}
\providecommand{\Mult}{\ensuremath{\mathrm{Mult}}}
\providecommand{\Add}{\ensuremath{\mathrm{Add}}}
\providecommand{\Exp}{\ensuremath{\mathrm{Exp}}}
\providecommand{\Zero}{\ensuremath{0}}
\providecommand{\Succ}{\ensuremath{\mathrm{S}}}
\providecommand{\lambd}{\ensuremath{\lambda}}
\providecommand{\bredpar}{\mathrel{\rightarrow_\beta}}
\providecommand{\bcd}[1]{\ensuremath{\ulcorner #1 \urcorner}}
\providecommand{\Set}[1]{\ensuremath{\{#1\}}}
\providecommand{\indcase}[1]{\par\noindent\emph{Case #1}.\quad}
\providecommand{\indfrm}[1]{\par\noindent\emph{Form #1}.\quad}
\providecommand{\iftag}[1]{#1}
\providecommand{\tagitem}[1]{\item[#1]}
\providecommand{\Elim}[1]{}
\providecommand{\Intro}[1]{}
\providecommand{\LeftR}[1]{}
\providecommand{\RightR}[1]{}
\providecommand{\TAss}[1]{}
\providecommand{\TRule}[1]{}

\makeatother
\providecommand{\olpath}{.}

% ===================== OLP compatibility layer (2) =====================
% --- second compatibility block: notation and shorthands used by the frozen OLP units ---
\providecommand{\CutCS}{\mathrm{Cut}}
\providecommand{\LeftSubproofLabel}{\relax}
\providecommand{\RightSubproofLabel}{\relax}
\providecommand{\maeh}{\relax}
\providecommand{\cheight}{\relax}
\providecommand{\ande}{\mathrm{and}_E}
\providecommand{\andi}{\mathrm{and}_I}
\providecommand{\ore}{\mathrm{or}_E}
\providecommand{\ori}{\mathrm{or}_I}
\providecommand{\typeof}{\mathord{:}}
\providecommand{\pValue}[1]{#1}
\providecommand{\Sent}{\mathrm{Sent}}
\providecommand{\cardplus}{\mathbin{+}}
\providecommand{\cardtimes}{\mathbin{\cdot}}
\providecommand{\cardexpo}{\mathbin{\hat{\ }}}
\providecommand{\cardnless}{\not<}
\providecommand{\cardle}{\leq}
\providecommand{\funfromto}[3]{#1\colon #2 \to #3}
\providecommand{\funimage}[1]{\operatorname{im}(#1)}
\providecommand{\fdefined}{\mathord{\downarrow}}
\providecommand{\canonord}{\operatorname{ord}}
\providecommand{\ordtimes}{\mathbin{\cdot}}
\providecommand{\ordeq}[2]{#1 = #2}
\providecommand{\ordneq}[2]{#1 \neq #2}
\providecommand{\ordtype}[1]{\operatorname{ord}(#1)}
\providecommand{\ordsucc}{\mathord{+1}}
\providecommand{\ordunion}{\cup}
\providecommand{\onesphere}{S^1}
\providecommand{\rotationsgroup}{\mathrm{SO}(2)}
\providecommand{\equivrep}[1]{[#1]}
\providecommand{\rlexless}{\prec}
\providecommand{\ZFC}{\mathrm{ZFC}}
\providecommand{\ZF}{\mathrm{ZF}}
\providecommand{\Rat}{\mathbb{Q}}
\providecommand{\Real}{\mathbb{R}}
\providecommand{\PosInt}{\mathbb{Z}^{+}}
\providecommand{\Id}{\mathrm{id}}
\providecommand{\Bin}{\mathrm{Bin}}
\providecommand{\pto}{\rightharpoonup}
\providecommand{\dom}{\operatorname{dom}}
\providecommand{\ran}{\operatorname{ran}}
\providecommand{\len}{\operatorname{len}}
\providecommand{\concat}{\mathbin{{}^\frown}}
\providecommand{\card}[1]{\lvert #1 \rvert}
\providecommand{\Pow}{\mathcal{P}}
\providecommand{\Setabs}[2]{\{\, #1 \mid #2 \,\}}
\providecommand{\comp}{\circ}
\providecommand{\Char}{\chi}
\providecommand{\Prf}{\operatorname{Prf}}
\providecommand{\Prov}{\operatorname{Prov}}
\providecommand{\Proj}{\operatorname{proj}}
\providecommand{\Gn}{\operatorname{gn}}
\providecommand{\gn}{\operatorname{gn}}
\providecommand{\Value}[1]{\operatorname{val}(#1)}
\providecommand{\Deduce}{\mathrm{D}}
\providecommand{\Dual}{\mathrm{d}}
\providecommand{\Ax}{\mathrm{Ax}}
\providecommand{\pSat}[3]{#1 \models #2\,[#3]}
\providecommand{\pAssign}{\sigma}
\providecommand{\bforall}{\forall}
\providecommand{\bexists}{\exists}
\providecommand{\liff}{\mathbin{\leftrightarrow}}
\providecommand{\FV}{\operatorname{FV}}
\providecommand{\UnaryInf}{\ensuremath{\mid}}
\providecommand{\BinaryInf}{\ensuremath{\mid}}
\providecommand{\depth}{\operatorname{depth}}
\providecommand{\oliflabeldef}[3]{#1}
\providecommand{\TMblank}{\sqcup}
\providecommand{\TMstroke}{\mid}
\providecommand{\TMstay}{\mathrm{S}}
\providecommand{\TMright}{\mathrm{R}}
\providecommand{\TMleft}{\mathrm{L}}
\providecommand{\TMendtape}{\triangleright}
\providecommand{\umin}{\mu}
\providecommand{\Sum}{\sum}
\providecommand{\citep}[1]{\cite{#1}}
\providecommand{\citet}[1]{\cite{#1}}
\providecommand{\citeauthor}[1]{\cite{#1}}
\providecommand{\citeyear}[1]{\cite{#1}}

% ===================== OLP compatibility layer (3) =====================
% --- third compatibility block: remaining notation seen in the frozen OLP units ---
\providecommand{\iso}{\cong}
\providecommand{\nszero}{\mathrm{0}}
\providecommand{\nssucc}{\mathrm{S}}
\providecommand{\nsplus}{\mathbin{+}}
\providecommand{\nssless}{<}
\providecommand{\nsless}{<}
\providecommand{\PL}{\mathrm{PL}}
\providecommand{\Nec}{\mathord{\Box}}
\providecommand{\RK}{\mathrm{RK}}
\providecommand{\ST}{\mathrm{ST}}
\providecommand{\SSubst}[3]{#1[#2/#3]}
\providecommand{\subst}[3]{#1[#2/#3]}
\providecommand{\tagfalse}{\mathrm{F}}
\providecommand{\cutr}{\mathrm{Cut}}
\providecommand{\cardneq}{\neq}
\providecommand{\issucc}[1]{\operatorname{succ}(#1)}
\providecommand{\tup}[1]{\langle #1 \rangle}
\providecommand{\Var}{\mathrm{Var}}
\providecommand{\Bew}{\operatorname{Bew}}
\providecommand{\Con}{\operatorname{Con}}
\providecommand{\Cons}{\operatorname{Cons}}
\providecommand{\pDeduce}{\mathrm{D}}
\providecommand{\psat}{\models}
\providecommand{\den}[1]{\llbracket #1 \rrbracket}
\providecommand{\Sem}[2]{\llbracket #1 \rrbracket_{#2}}
\providecommand{\DomainOf}[1]{\lvert #1 \rvert}
\providecommand{\Form}{\mathrm{Form}}
\providecommand{\Terms}{\mathrm{Term}}
\providecommand{\Atom}{\mathrm{Atom}}
\providecommand{\FVars}[1]{\operatorname{FV}(#1)}
\providecommand{\fv}[1]{\operatorname{fv}(#1)}
\providecommand{\tuplep}[1]{\langle #1 \rangle}
\providecommand{\restrict}[2]{#1 \upharpoonright #2}
\providecommand{\sto}{\rightarrow}
\providecommand{\FmlaSet}[1]{\mathrm{Fmla}_{#1}}
\providecommand{\Thm}{\mathrm{Thm}}
\providecommand{\Eval}{\operatorname{val}}
\providecommand{\blank}{\sqcup}
\providecommand{\blankTape}{\sqcup}
\providecommand{\Consis}{\operatorname{Con}}
\providecommand{\drs}{\mathrm{d}}
\providecommand{\VRule}{\mid}
\providecommand{\Bet}{\mathrm{B}}
\providecommand{\Seg}[1]{\mathrm{Seg}_{#1}}
\providecommand{\Enums}{\mathrm{En}}
\providecommand{\Z}{\mathbb{Z}}
\providecommand{\N}{\mathbb{N}}
\providecommand{\abs}[1]{\lvert #1 \rvert}
\providecommand{\floor}[1]{\lfloor #1 \rfloor}
\providecommand{\ceil}[1]{\lceil #1 \rceil}

% ===================== OLP compatibility layer (4) =====================
% --- fourth compatibility block: harvested from the per-candidate compile logs ---
\providecommand{\beredpar}{\mathrel{\rightarrow_\beta}}
\makeatletter
\def\bered{\@ifnextchar[{\bered@i}{\mathrel{\rightarrow_\beta}}}
\def\bered@i[#1]{\mathrel{\rightarrow_\beta}}
\def\bred{\@ifnextchar[{\bred@i}{\mathrel{\rightarrow_\beta}}}
\def\bred@i[#1]{\mathrel{\rightarrow_\beta}}
\makeatother
\providecommand{\beredone}{\mathrel{=_\beta}}
\providecommand{\bredone}{\mathrel{=_\beta}}
\providecommand{\xredone}{\mathrel{=_\beta}}
\providecommand{\aconvone}{\mathrel{=_\alpha}}
\providecommand{\cardless}{<}
\providecommand{\tagrefs}[1]{#1}
\providecommand{\ordplus}{\mathbin{+}}
\providecommand{\funrestrictionto}[2]{#1 \upharpoonright #2}
\providecommand{\fundefined}{\mathord{\uparrow}}
\providecommand{\Trm}{\mathrm{Trm}}
\providecommand{\OProv}{\operatorname{Prov}}
\providecommand{\OPrf}{\operatorname{Prf}}
\providecommand{\aeq}{\mathrel{=}}
\providecommand{\mClass}[1]{[#1]}
\providecommand{\equivclass}[1]{[#1]}
\providecommand{\pheight}{\operatorname{ht}}
\providecommand{\phy}{\operatorname{ht}}
\providecommand{\Cut}{\mathrm{Cut}}
\providecommand{\QuantRank}{\operatorname{qr}}
\providecommand{\elemequiv}{\mathrel{\equiv}}
\providecommand{\PIso}{\mathrm{PIso}}
\providecommand{\Theory}{\mathrm{Th}}
\providecommand{\supstrict}{\subsetneq}
\providecommand{\trcl}{\operatorname{trcl}}
\providecommand{\cref}[1]{\ref{#1}}
\providecommand{\Discharge}{\mathrm{Dis}}
\providecommand{\Weakening}{\mathrm{Wk}}
\providecommand{\ZFminus}{\mathrm{ZF}^{-}}
\makeatletter
\renewcommand{\Dual}{\@ifnextchar[{\Dual@i}{\ifmmode\mathrm{d}\else$\mathrm{d}$\fi}}
\def\Dual@i[#1]{\ifmmode\mathrm{d}[#1]\else$\mathrm{d}[#1]$\fi}
\makeatother
\providecommand{\RightR}{\mathrm{R}}
\providecommand{\LeftR}{\mathrm{L}}
\providecommand{\Intro}{\mathrm{I}}
\providecommand{\Axiom}{\mathrm{Ax}}
\providecommand{\formula}[1]{#1}
\providecommand{\stared}{\ast}
\providecommand{\oset}[2]{\{\, #1 \mid #2 \,\}}
\providecommand{\setof}[2]{\{\, #1 \mid #2 \,\}}
\providecommand{\pow}[1]{\mathcal{P}(#1)}
\providecommand{\ordpair}[2]{\langle #1, #2 \rangle}
\providecommand{\ordtriple}[3]{\langle #1, #2, #3 \rangle}
\providecommand{\union}{\cup}
\providecommand{\inter}{\cap}
\providecommand{\seq}[1]{\langle #1 \rangle}
\providecommand{\SuchThat}{\mid}
\providecommand{\where}{\mid}
\providecommand{\TVar}{\mathrm{TVar}}
\providecommand{\Bool}{\mathrm{Bool}}
\providecommand{\Val}{\operatorname{Val}}
\providecommand{\TL}[1]{\mathrm{TL}}
\providecommand{\xstared}{\ast}

% ===================== OLP compatibility layer (5) =====================
% --- fifth compatibility block: rescues the quarantined OLP candidates ---
\providecommand{\ident}{\operatorname{id}}
\providecommand{\setrank}{\operatorname{rank}}
\providecommand{\cardfont}{\mathfrak{c}}
\providecommand{\Complement}{\overline}
\providecommand{\QR}{\mathrm{QR}}
\providecommand{\becd}{\ensuremath{\ulcorner}}
\providecommand{\disjointsum}{\mathbin{\oplus}}
\providecommand{\cardeq}[2]{#1 = #2}
\providecommand{\xred}{\mathrel{\rightarrow}}
% --- sixth compatibility block: names reported by real compile logs ---
\providecommand{\ORefut}{\operatorname{Refut}}
\providecommand{\ORProv}{\operatorname{Prov}}
\providecommand{\cutrank}{\operatorname{cr}}
\providecommand{\Refut}{\operatorname{Refut}}
\providecommand{\DischargeRule}{\mathrm{Dis}}
\providecommand{\Contraction}{\mathrm{C}}
\providecommand{\bottomAlignProof}{\relax}
\providecommand{\FalseCl}{\bot}
\providecommand{\red}{\mathrel{\rightarrow}}
\providecommand{\varAssign}{\sigma}
\providecommand{\defiff}{\mathrel{\leftrightarrow}}
\providecommand{\FalseInt}{\bot}
\providecommand{\fn}{\operatorname{fn}}
\providecommand{\dotted}{\cdot}

% --- tag-macro semantic fix (2026-10-13, final): \iftag CHOOSES a branch and the OLP proof body is
% always the SECOND group (\iftag{tag}{body}, {tag}{body}{c}, {tag}{body}{c}{d}); a definition that
% prints the tag literal drops the proof body. Bodies contain blank lines, hence \long.
\makeatletter
\providecommand{\iftag}{}
\long\def\iftag#1#2{\@ifnextchar\bgroup{\iftag@three{#2}}{\iftag@stop{#2}}}
\long\def\iftag@three#1#2{\iftag@stop{#1}\@ifnextchar\bgroup{\iftag@eat}{\relax}}
\long\def\iftag@eat#1{}
\long\def\iftag@stop#1{#1}
\long\def\oliflabeldef#1#2{\@ifnextchar\bgroup{\@gobble}{\iftag@stop{#2}}}
\providecommand{\tagrefs}[1]{}
\providecommand{\settag}[1]{}
\makeatother
\begin{document}

% ==== OLP file: content/incompleteness/representability-in-q/basic-representable.tex ====
% Part: incompleteness
% Chapter: representability-in-q
% Section: basic-representable



\olfileid{inc}{req}{bre}
\olsection{Basic Functions are Representable in~$\Th{Q}$}

First we have to show that all the basic functions are representable
in~$\Th{Q}$. In the end, we need to show how to assign to each $k$-ary
basic function $f(x_0,\dots,x_{k-1})$ !!a{formula}
$!A_f(x_0,\dots,x_{k-1},y)$ that represents it.

We will be able to represent zero, successor, plus, times, the
characteristic function for equality, and projections. In each case,
the appropriate representing function is entirely straightforward; for
example, zero is represented by the formula $y = \Obj 0$, successor is
represented by the !!{formula} $x_0' = y$, and addition is represented
by the !!{formula} $(x_0 + x_1) = y$. The work involves showing that
$\Th{Q}$ can prove the relevant !!{sentence}s; for example, saying
that addition is represented by the !!{formula} above involves showing
that for every pair of natural numbers $m$ and $n$, $\Th{Q}$ proves
\begin{align*}
& \eq[\num n + \num m][\num {n+m}] \text{ and}\\
& \lforall[y][(\eq[(\num n + \num m)][y] \lif \eq[y][\num{n+m}])].
\end{align*}

\begin{prop}
\ollabel{prop:rep-zero}
The zero function $\Zero(x) = 0$ is represented in~$\Th{Q}$ by
$!A_{\Zero}(x,y) \ident \eq[y][\Obj 0]$.
\end{prop}

\begin{prop}
\ollabel{prop:rep-succ}
The successor function $\Succ(x) = x+1$ is represented in~$\Th{Q}$ by
$!A_{\Succ}(x,y) \ident \eq[y][x']$.  
\end{prop}

\begin{prop}
\ollabel{prop:rep-proj}
The projection function $\Proj{n}{i}(x_0, \dots, x_{n-1}) = x_i$ is
represented in~$\Th{Q}$ by \[!A_{\Proj{n}{i}}(x_0, \dots, x_{n-1}, y) \ident \eq[y][x_i].\]
\end{prop}

\begin{prob}
Prove that $\eq[y][\Obj 0]$, $\eq[y][x']$, and $\eq[y][x_i]$ represent
$\Zero$, $\Succ$, and $\Proj{n}{i}$, respectively.
\end{prob}

\begin{prop}
\ollabel{prop:rep-id}
The characteristic function of~$=$,
\[
\Char{=}(x_0, x_1) =
\begin{cases}
  1 & \text{if } x_0 =x_1\\
  0 & otherwise
\end{cases}
\]
is represented in~$\Th{Q}$ by
\[
  !A_{\Char{=}}(x_0, x_1, y) \ident (\eq[x_0][x_1] \land \eq[y][\num{1}]) \lor (\eq/[x_0][x_1] \land
\eq[y][\num{0}]).
\]
\end{prop}

The proof requires the following lemma.

\begin{lem}
\ollabel{lem:q-proves-neq} Given natural numbers $n$ and $m$, if $n
\neq m$, then $\Th{Q} \Proves \eq/[\num n][\num m]$.
\end{lem}

\begin{proof}
Use induction on $n$ to show that for every $m$, if $n \neq m$, then
$Q \Proves \eq/[\num n][\num m]$.

In the base case, $n = 0$. If $m$ is not equal to $0$, then $m = k +
1$ for some natural number $k$. We have an axiom that says
$\lforall[x][\eq/[0][x']]$. By a quantifier axiom, replacing $x$ by
  $\num k$, we can conclude $\eq/[0][\num k']$. But $\num k'$ is just
  $\num m$.

In the induction step, we can assume the claim is true for $n$, and
consider $n+1$. Let $m$ be any natural number. There are two
possibilities: either $m = 0$ or for some $k$ we have $m = k+1$. The
first case is handled as above. In the second case, suppose $n+1 \neq
k+1$. Then $n \neq k$. By the induction hypothesis for $n$ we have
$\Th{Q} \Proves \eq/[\num n][\num k]$. We have an axiom that says
$\lforall[x][\lforall[y][\eq[x'][y'] \lif \eq[x][y]]]$. Using a
quantifier axiom, we have $\eq[\num n'][\num k'] \lif \eq[\num n][\num
  k]$. Using propositional logic, we can conclude, in $\Th{Q}$,
$\eq/[\num n][\num k] \lif \eq/[\num n'][\num k']$. Using modus
ponens, we can conclude $\eq/[\num n'][\num k']$, which is what we want,
since $\num k'$ is $\num m$.
\end{proof}

\begin{explain}
Note that the lemma does not say much: in essence it says that $\Th{Q}$ can
prove that different numerals denote different objects. For example,
$\Th{Q}$ proves $0'' \neq 0'''$. But showing that this holds in general
requires some care. Note also that although we are using induction, it
is induction \emph{outside} of $\Th{Q}$.
\end{explain}

\begin{proof}[Proof of \olref{prop:rep-id}]
If $n = m$, then $\num{n}$ and $\num{m}$ are the same term, and
$\Char{=}(n, m) = 1$. But $\Th{Q} \Proves (\eq[\num{n}][\num{m}] \land
\eq[\num{1}][\num{1}])$, so it proves $!A_=(\num{n}, \num{m},
\num{1})$.  If $n \neq m$, then $\Char=(n, m) = 0$. By
\olref{lem:q-proves-neq}, $\Th{Q} \Proves \eq/[\num{n}][\num{m}]$ and
so also $(\eq/[\num{n}][\num{m}] \land \Obj 0 = \Obj 0)$. Thus $\Th{Q}
\Proves !A_=(\num{n}, \num{m}, \num{0})$.

For the second part, we also have two cases. If $n = m$, we have to
show that $\Th{Q} \Proves \lforall[y][(!A_=(\num{n}, \num{m}, y)
  \lif \eq[y][\num{1}])]$.  Arguing informally, suppose $!A_=(\num{n},
\num{m}, y)$, i.e.,
\[
(\eq[\num{n}][\num{n}] \land \eq[y][\num{1}]) \lor
(\eq/[\num{n}][\num{n}] \land \eq[y][\num{0}])
\]
The left disjunct implies $\eq[y][\num{1}]$ by logic; the right
contradicts $\eq[\num{n}][\num{n}]$ which is provable by logic.

Suppose, on the other hand, that $n \neq m$. Then $!A_=(\num{n},
\num{m}, y)$ is
\[
(\eq[\num{n}][\num{m}] \land \eq[y][\num{1}]) \lor
(\eq/[\num{n}][\num{m}] \land \eq[y][\num{0}])
\]
Here, the left disjunct contradicts $\eq/[\num{n}][\num{m}]$, which is
provable in $\Th{Q}$ by \olref{lem:q-proves-neq}; the right disjunct
entails $\eq[y][\num{0}]$.
\end{proof}

\begin{prop}
\ollabel{prop:rep-add}
The addition function $\Add(x_0, x_1) = x_0+x_1$ is represented
in~$\Th{Q}$ by
\[
  !A_{\Add}(x_0, x_1, y) \ident \eq[y][(x_0 + x_1)].
\]
\end{prop}

\begin{lem}
\ollabel{lem:q-proves-add}
$\Th{Q} \Proves \eq[(\num{n} + \num{m})][\num{n+m}]$
\end{lem}

\begin{proof}
We prove this by induction on~$m$. If $m = 0$, the claim is that
$\Th{Q} \Proves \eq[(\num{n} + \Obj 0)][\num{n}]$. This follows by
axiom~$!Q_4$.  Now suppose the claim for $m$; let's prove the claim
for $m+1$, i.e., prove that $\Th{Q} \Proves \eq[(\num{n} +
  \num{m+1})][\num{n+m+1}]$. Note that $\num{m+1}$ is just $\num{m}'$,
and $\num{n+m+1}$ is just $\num{n+m}'$.  By axiom $!Q_5$, $\Th{Q}
\Proves \eq[(\num{n} + \num{m}')][(\num{n}+\num{m})']$. By induction
hypothesis, $\Th{Q} \Proves \eq[(\num{n} + \num{m})][\num{n+m}]$. So
$\Th{Q} \Proves \eq[(\num{n} + \num{m}')][\num{n+m}']$.
\end{proof}

\begin{proof}[Proof of \olref{prop:rep-add}]
The !!{formula}~$!A_\Add(x_0, x_1, y)$ representing $\Add$ is
$\eq[y][(x_0 + x_1)]$. First we show that if $\Add(n, m) = k$, then
$\Th{Q} \Proves !A_\Add(\num{n}, \num{m}, \num{k})$, i.e., $\Th{Q}
\Proves \eq[\num{k}][(\num{n} + \num{m})]$. But since $k = n + m$,
$\num{k}$ just is $\num{n+m}$, and we've shown in
\olref{lem:q-proves-add} that $\Th{Q} \Proves \eq[(\num{n} +
  \num{m})][\num{n+m}]$.

We also have to show that if $\Add(n, m) = k$, then
\[
\Th{Q} \Proves \lforall[y][(!A_\Add(\num{n}, \num{m}, y) \lif
  \eq[y][\num{k}])].
\]
Suppose we have $\eq[(\num{n} + \num{m})][y]$. Since
\[
\Th{Q} \Proves \eq[(\num{n}+\num{m})][\num{n+m}],
\]
we can replace the left side with $\num{n+m}$ and get
$\eq[\num{n+m}][y]$, for arbitrary~$y$.
\end{proof}

\begin{prop}
\ollabel{prop:rep-mult}
The multiplication function $\Mult(x_0, x_1) = x_0 \cdot x_1$ is represented
in~$\Th{Q}$ by
\[
  !A_{\Mult}(x_0, x_1, y) \ident y = (x_0 \times x_1).
\]
\end{prop}

\begin{proof} Exercise. \end{proof}

\begin{lem}
\ollabel{lem:q-proves-mult}
$\Th{Q} \Proves \eq[(\num{n} \times \num{m})][\num{n \cdot m}]$
\end{lem}

\begin{proof} Exercise. \end{proof}

\begin{prob}
Prove \olref[inc][req][bre]{lem:q-proves-mult}.
\end{prob}

\begin{prob}
Use \olref[inc][req][bre]{lem:q-proves-mult} to prove
\olref[inc][req][bre]{prop:rep-mult}.
\end{prob}

\begin{explain}
  Recall that we use $\times$ for the function symbol of the language
  of arithmetic, and $\cdot$ for the ordinary multiplication operation
  on numbers. So $\cdot$ can appear between expressions for numbers
  (such as in $m \cdot n$) while $\times$ appears only between terms
  of the language of arithmetic (such as in $(\num{m} \times
  \num{n})$).  Even more confusingly, $+$ is used for both the
  !!{function} and the addition operation. When it appears between
  terms---e.g., in $(\num{n} + \num{m})$---it is the $2$-place
  !!{function} of the language of arithmetic, and when it appears
  between numbers---e.g., in $n+m$---it is the addition operation.
  This includes the case $\num{n+m}$: this is the standard numeral
  corresponding to the number~$n+m$.
\end{explain}


% ==== OLP file: content/incompleteness/representability-in-q/beta-function.tex ====
% Part: incompleteness
% Chapter: representability-in-q
% Section: beta-function



\olfileid{inc}{req}{bet}
\olsection{The Beta Function Lemma}


In order to show that we can carry out primitive recursion if
addition, multiplication, and $\Char{=}$ are available, we need to
develop functions that handle sequences. (If we had exponentiation as
well, our task would be easier.) When we had primitive recursion, we
could define things like the ``$n$-th prime,'' and pick a fairly
straightforward coding. But here we do not have primitive
recursion---in fact we want to show that we can do primitive recursion
using minimization---so we need to be more clever.

\begin{lem}
\ollabel{lem:beta}
There is a function $\beta(d,i)$ such that for every sequence $a_0$,
\dots,~$a_n$ there is a number~$d$, such that for every $i \le n$,
$\beta(d,i) = a_i$. Moreover, $\beta$ can be defined from the basic
functions using just composition and regular minimization.
\end{lem}

Think of $d$ as coding the sequence $\tuple{a_0, \dots, a_n}$, and
$\beta(d,i)$ returning the $i$-th element. (Note that this ``coding''
does \emph{not} use the power-of-primes coding we're already familiar
with!). The lemma is fairly minimal; it doesn't say we can concatenate
sequences or append elements, or even that we can \emph{compute}~$d$
from $a_0$, \dots,~$a_n$ using functions definable by composition and
regular minimization. All it says is that there is a ``decoding''
function such that every sequence is ``coded.''

The use of the notation $\beta$ is G\"odel's. To repeat, the hard part
of proving the lemma is defining a suitable~$\beta$ using the
seemingly restricted resources, i.e., using just composition and
minimization---however, we're allowed to use addition, multiplication,
and~$\Char{=}$. There are various ways to prove this lemma, but one of
the cleanest is still G\"odel's original method, which used a
number-theoretic fact called Sunzi's Theorem
(traditionally, the ``Chinese Remainder Theorem'').

\begin{defn}
Two natural numbers $a$ and $b$ are \emph{relatively prime} iff their
greatest common divisor is~$1$; in other words, they have no other
divisors in common.
\end{defn}

\begin{defn}
Natural numbers $a$ and $b$ are \emph{congruent modulo~$c$},
$a \equiv b \mod c$, iff $c \mid (a-b)$, i.e., $a$ and $b$ have the
same remainder when divided by~$c$.
\end{defn}

Here is Sunzi's Theorem:
\begin{thm}
Suppose $x_0$, \dots,~$x_n$ are (pairwise) relatively prime. Let
$y_0$, \dots,~$y_n$ be any numbers. Then there is a number $z$ such that
\begin{align*}
z & \equiv y_0 \mod x_0 \\
z & \equiv y_1 \mod x_1 \\
& \vdots  \\
z & \equiv y_n \mod x_n.
\end{align*}
\end{thm}

Here is how we will use Sunzi's Theorem: if $x_0$,
\dots,~$x_n$ are bigger than $y_0$, \dots,~$y_n$ respectively, then we
can take $z$ to code the sequence $\tuple{y_0, \dots,y_n}$. To
recover~$y_i$, we need only divide $z$ by~$x_i$ and take the
remainder. To use this coding, we will need to find suitable values
for $x_0$, \dots,~$x_n$.

A couple of observations will help us in this regard. Given
$y_0$, \dots,~$y_n$, let
\begin{align*}
j &= \max(n, y_0 + 1, \dots, y_n + 1), \\
m &= \lcm(1,\dots,j),
\end{align*}
and let
\begin{align*}
x_0 & = 1 + m \\
x_1 & = 1 + 2 \cdot m \\
x_2 & = 1 + 3 \cdot m \\
& \vdots  \\
x_n & = 1 + (n+1) \cdot m
\end{align*}
Then two things are true:
\begin{enumerate}
\item\ollabel{rel-prime} $x_0,\dots,x_n$ are relatively prime.
\item\ollabel{less} For each $i$, $y_i < x_i$.
\end{enumerate}
To see that \olref{rel-prime} is true, note that if $p$ is a prime number
and $p \mid x_i$ and $p \mid x_k$, then $p \mid 1 + (i+1) m$ and
$p \mid 1 + (k+1) m$. But then $p$ divides their difference,
\[
(1 + (i+1)m) - (1+ (k+1)m) = (i-k) m.
\]
Since $p$ divides $1 + (i+1)m$, it can't divide $m$ as well
(otherwise, the first division would leave a remainder of~$1$). So $p$
divides $i-k$, since $p$ divides $(i-k)m$. But $\left|i-k\right|$ is at
most~$n$, and we have chosen $j \geq n$, so this implies that
$p \mid m$, again a contradiction. So there is no prime number dividing
both $x_i$ and $x_k$. Clause~\olref{less} is easy:
we have $y_i < j \leq m < x_i$.

Now let us prove the $\beta$ function lemma. Remember that we can use
$0$, successor, plus, times, $\Char{=}$, projections, and any function
defined from them using composition and minimization applied to
regular functions. We can also use a relation if its characteristic
function is so definable. As before we can show that these relations
are closed under Boolean combinations and bounded quantification; for
example:
\begin{align*}
\fn{not}(x) & \defis \Char{=}(x,0)\\
\bmin{x \leq z}{R(x,y)} & \defis \umin{x}{(R(x,y) \lor x = z)}\\
\bexists{x \leq z}{R(x,y)} & \defiff R(\bmin{x \leq z}{R(x,y)}, y)
\end{align*}
We can then show that all of the following are also definable without
primitive recursion:
\begin{enumerate}
\item The pairing function, $J(x,y) = \frac{1}{2}[(x+y)(x+y+1)] + x$;
% maybe explain more what is going on here, a bit confusing.
\item the projection functions
\begin{align*}
K(z) & = \bmin{x \leq z}{\bexists{y \leq z}{z = J(x,y)}},\\
L(z) & = \bmin{y \leq z}{\bexists{x \leq z}{z = J(x,y)}};
\end{align*}
\item the less-than relation $x < y$;
\item the divisibility relation $x \mid y$;
% \item $x \tsub y$
% \item $\fn{Prime}(x)$
% \item Assuming $p$ is prime, the relation ``$x$ is a power of $p$'':
% \[
% \bforall{y \leq x}{(y \mid x \lif y = 1 \lor y = x)}.
% \]
\item the function $\fn{rem}(x,y)$ which returns the remainder when
  $y$ is divided by~$x$.
\end{enumerate}
Now define
\begin{align*}
\beta^*(d_0,d_1,i) & = \fn{rem}(1+(i+1) d_1,d_0) \text{ and}\\
\beta(d,i) & = \beta^*(K(d),L(d),i).
\end{align*}
This is the function we want. Given $a_0,\dots,a_n$ as above, let
\[
j = \max(n,a_0+1,\dots,a_n+1),
\]
and let $d_1 = \lcm(1,\dots,j)$. By \olref{rel-prime} above,
we know that $1+d_1$, $1+2 d_1$, \dots, $1+(n+1) d_1$ are relatively
prime, and by~\olref{less} that all are greater than $a_0,\dots,a_n$.
By Sunzi's Theorem there is a value~$d_0$ such that for each~$i$,
\[
d_0 \equiv a_i \mod (1+(i+1)d_1)
\]
and so (because $d_1$ is greater than~$a_i$),
\[
a_i = \fn{rem}(1+(i+1)d_1,d_0).
\]
Let $d = J(d_0,d_1)$. Then for each $i \le n$, we have
\begin{align*}
\beta(d,i) & =  \beta^*(d_0,d_1,i) \\
& =  \fn{rem}(1+(i+1) d_1,d_0) \\
& =  a_i
\end{align*}
which is what we need. This completes the proof of the
$\beta$-function lemma.

\begin{prob}
  Show that the relations $x < y$, $x \mid y$, and the
  function~$\fn{rem}(x,y)$ can be defined without primitive recursion.
  You may use $0$, successor, plus, times, $\Char{=}$, projections,
  and bounded minimization and quantification.
\end{prob}


% ==== OLP file: content/incompleteness/representability-in-q/c-representable.tex ====
% Part: incompleteness
% Chapter: representability-in-q
% Section: c-representable



\olfileid{inc}{req}{cre}
\olsection{Functions in $C$ are Representable in $\Th{Q}$}

We have to show that every function in $C$ is representable in
$\Th{Q}$. In the end, we need to show how to assign to each $k$-ary
function $f(x_0,\dots,x_{k-1})$ in $C$ a formula
$!A_f(x_0,\dots,x_{k-1},y)$ that represents it.

\begin{lem}
Given natural numbers $n$ and $m$, if $n \neq m$, then $Q \Proves \num
n \neq \num m$.
\end{lem}

\begin{proof}
Use induction on $n$ to show that for every $m$, if $n \neq m$, then
$Q \Proves \eq/[\num n][\num m]$.

In the base case, $n = 0$. If $m$ is not equal to $0$, then $m = k +
1$ for some natural number $k$. We have an axiom that says
$\lforall[x][\eq/[0][x']]$. By a quantifier axiom, replacing $x$ by
  $\num k$, we can conclude $\eq/[0][\num k']$. But $\num k'$ is just
  $\num m$.

In the induction step, we can assume the claim is true for $n$, and
consider $n+1$. Let $m$ be any natural number. There are two
possibilities: either $m = 0$ or for some $k$ we have $m = k+1$. The
first case is handled as above. In the second case, suppose $n+1 \neq
k+1$. Then $n \neq k$. By the induction hypothesis for $n$ we have
$\Th{Q} \Proves \eq/[\num n][\num k]$. We have an axiom that says
$\lforall[x][\lforall[y][\eq[x'][y'] \lif \eq[x][y]]]$. Using a
quantifier axiom, we have $\eq[\num n'][\num k'] \lif \eq[\num n][\num
  k]$. Using propositional logic, we can conclude, in $\Th{Q}$,
$\eq/[\num n][\num k] \lif \eq/[\num n'][\num k']$. Using modus
ponens, we can conclude $\eq/[\num n'][\num k']$, which is what we want,
since $\num k'$ is $\num m$.
\end{proof}

Note that the lemma does not say much: in essence it says that $\Th{Q}$ can
prove that different numerals denote different objects. For example,
$\Th{Q}$ proves $0'' \neq 0'''$. But showing that this holds in general
requires some care. Note also that although we are using induction, it
is induction \emph{outside} of $\Th{Q}$.

We will be able to represent zero, successor, plus, times, the
characteristic function for equality, and projections. In each case,
the appropriate representing function is entirely straightforward; for
example, zero is represented by the formula
\[
y = 0,
\]
successor is represented by the formula
\[
x_0' = y,
\]
and plus is represented by the formula
\[
x_0 + x_1 = y.
\]
The work involves showing that $\Th{Q}$ can prove the relevant statements;
for example, saying that plus is represented by the formula above
involves showing that for every pair of natural numbers $m$ and $n$,
$\Th{Q}$ proves
\[
\eq[\num n + \num m][\num {n+m}]
\]
and
\[
\lforall[y][(\eq[(\num n + \num m)][y] \lif \eq[y][\num {n+m}])].
\]

What about composition? Suppose $h$ is defined by
\[
h(x_0,\dots,x_{l-1}) = f(g_0(x_0,\dots,x_{l-1}), \dots,
g_{k-1}(x_0,\dots,x_{l-1})).
\]
where we have already found formulas $!A_f, !A_{g_0}, \dots,
!A_{g_{k-1}}$ representing the functions $f,g_0,\dots,g_{k-1}$,
respectively. Then we can define a formula $!A_h$ representing $h$,
by defining $!A_h(x_0,\dots,x_{l-1},y)$ to be
\begin{multline*}
  \lexists[z_0\dots][\lexists[z_{k-1}][(!A_{g_0}(x_0,\dots,x_{l-1},z_0) \land
  \dots \land !A_{g_{k-1}}(x_0,\dots,x_{l-1},z_{k-1}) \land]]\\
  !A_f(z_0,\dots,z_{k-1},y)).
\end{multline*}

Finally, let us consider unbounded search. Suppose $g(x,\vec z)$ is
regular and representable in $\Th{Q}$, say by the formula $!A_g(x,\vec
z,y)$. Let $f$ be defined by $f(\vec z) = \mu x \; g(x,\vec z)$. We would
like to find a formula $!A_f(\vec z,y)$ representing $f$. Here is
a natural choice:
\[
!A_f(\vec z,y) \ident !A_g(y,\vec z,0) \land \lforall[w][(w < y
\lif \lnot !A_g(w,\vec z,0))].
\]

It can be shown that this works using some additional lemmas, e.g.,

\begin{lem}
  For every variable $x$ and every natural number $n$, $\Th{Q}$ proves $\eq[(x'
  + \num n)][(x + \num n)']$.
\end{lem}

It is again worth mentioning that this is weaker than saying that $\Th{Q}$
proves $\lforall[x][\lforall[y][(x' + y) = (x +y)']]$ (which is false).

\begin{proof}
The proof is, as usual, by induction on $n$. In the base case, $n =
0$, we need to show that $\Th{Q}$ proves $\eq[(x'+0)][(x+0)']$. But we have:
\begin{eqnarray*}
& & \eq[(x' + 0)][x'] \quad \text{from axiom 4} \\
& & \eq[(x + 0)][x] \quad \text{from axiom 4} \\
& & \eq[(x+0)'][x'] \quad \text{by line 2} \\
& & \eq[(x' + 0)][(x + 0)'] \quad \text{lines 1 and 3}
\end{eqnarray*}
In the induction step, we can assume that we have derived $\eq[(x' +
  \num n)][(x + \num n)']$ in $\Th{Q}$. Since $\num{n+1}$ is $\num
n'$, we need to show that $\Th{Q}$ proves $\eq[(x' + \num n')][(x + \num
n')']$. The following chain of equalities can be derived in $\Th{Q}$:
\begin{eqnarray*}
(x' + \num n') & \eq & (x' + \num n)' \quad \text{axiom 5} \\
& \eq & (x + \num n')' \quad \text{from the inductive hypothesis}
\end{eqnarray*}
\end{proof}

\begin{lem}
\begin{enumerate}
\item $\Th{Q}$ proves $\lnot (x < \num 0)$.
\item For every natural number $n$, $\Th{Q}$ proves
\[
x < \num {n+1} \lif (\eq[x][\num 0] \lor \dots \lor \eq[x][\num n]).
\]
\end{enumerate}
\end{lem}

\begin{proof}
Let us do 1 and part of 2, informally (i.e., only giving hints as to
how to construct the formal !!{derivation}).

For part 1, by the definition of $<$, we need to prove $\lnot
\lexists[y][\eq[(y'+x)][0]]$ in $\Th{Q}$, which is equivalent (using
the axioms and rules of first-order logic) to $\lforall[y][\eq/[(y' +
    x)][0]]$. Here is the idea: suppose $\eq[(y'+x)][0]$. If $x$ is 0,
we have $\eq[(y' + 0)][0]$. But by axiom 4 of $\Th{Q}$, we have
$\eq[(y' + 0)][y']$, and by axiom 2 we have $\eq/[y'][0]$, a
contradiction. So $\lforall[y][\eq/[(y' +x)][0]]$. If $x$ is not $0$,
by axiom 3 there is a $z$ such that $\eq[x][z']$. But then we have
$\eq[(y' + z')][0]$. By axiom~5, we have $\eq[(y' + z)'][0]$, again
contradicting axiom~2.

For part 2, use induction on $n$. Let us consider the base case, when
$n =0$. In that case, we need to show $x < \num 1 \lif \eq[x][\num
  0]$. Suppose $x < \num 1$. Then by the defining axiom for $<$, we
have $\lexists[y][\eq[(y'+x)][0']]$. Suppose $y$ has that property; so
we have $\eq[y'+x][0']$.

We need to show $\eq[x][0]$. By axiom 3, if $x$ is not 0, it is equal
to $z'$ for some $z$. Then we have $\eq[(y' + z')][0']$. By axiom~5 of
  $\Th{Q}$, we have $\eq[(y'+z)'][0']$. By axiom 1, we have $\eq[(y' +
    z)][0]$. But this means, by definition, $z < 0$, contradicting
  part~1.
\end{proof}

\begin{explain}
We have shown that the set of computable functions can be
characterized as the set of functions representable in $\Th{Q}$. In fact,
the proof is more general. From the definition of representability, it
is not hard to see that any theory extending $\Th{Q}$ (or in which one can
interpret $\Th{Q}$) can represent the computable functions; but,
conversely, in any !!{derivation} system in which the notion of proof is
computable, every representable function is computable. So, for
example, the set of computable functions can be characterized as the
set of functions represented in Peano arithmetic, or even Zermelo
Fraenkel set theory. As G\"odel noted, this is somewhat surprising.
We will see that when it comes to provability, questions are very
sensitive to which theory you consider; roughly, the stronger the
axioms, the more you can prove. But across a wide range of axiomatic
theories, the representable functions are exactly the computable ones.
\end{explain}


% ==== OLP file: content/incompleteness/representability-in-q/c.tex ====
% Part: incompleteness
% Chapter: representability-in-q
% Section: c



\olfileid{inc}{req}{cee}
\olsection{The Functions $C$}

Let $C$ be the smallest set of functions containing
\begin{enumerate}
\item $0$,
\item successor,
\item addition,
\item multiplication,
\item projections, and
\item the characteristic function for equality, $\Char{=}$;
\end{enumerate}
and closed under
\begin{enumerate}
\item composition, and
\item unbounded search, applied to regular
functions.
\end{enumerate}
Remember this last restriction means simply that you can only use the
$\mu$ operation when the result is total. Compare this to the
definition of the \emph{general recursive} functions: here we have
added plus, times, and $\Char{=}$, but we have dropped primitive
recursion.

Clearly everything in $C$ is recursive, since plus, times,
and $\Char{=}$ are. We will show that the converse is also true; this
amounts to saying that with the other stuff in $C$ we can carry out
primitive recursion.


% ==== OLP file: content/incompleteness/representability-in-q/comp-representable.tex ====
% Part: incompleteness
% Chapter: representability-in-q
% Section: comp-representable



\olfileid{inc}{req}{crq}
\olsection{Computable Functions are Representable in $\Th{Q}$}

\begin{thm}
Every computable function is representable in~$\Th{Q}$.
\end{thm}

\begin{proof}
For definiteness, and using the Church--Turing Thesis, let's say that a
function is computable iff it is general recursive. The general
recursive functions are those which can be defined from the zero
function~$\Zero$, the successor function~$\Succ$, and the projection
function~$\Proj{n}{i}$ using composition, primitive recursion, and
regular minimization. By \olref[pri]{lem:prim-rec}, any function~$h$
that can be defined from $f$ and~$g$ can also be defined using
composition and regular minimization from $f$, $g$, and $\Zero$,
$\Succ$, $\Proj{n}{i}$, $\Add$, $\Mult$, $\Char{=}$. Consequently, a
function is general recursive iff it can be defined from $\Zero$,
$\Succ$, $\Proj{n}{i}$, $\Add$, $\Mult$, $\Char{=}$ using composition
and regular minimization.

We've furthermore shown that the basic functions in question are
representable in~$\Th{Q}$
(\cref{inc:req:bre:prop:rep-zero,inc:req:bre:prop:rep-succ,inc:req:bre:prop:rep-proj,inc:req:bre:prop:rep-id,inc:req:bre:prop:rep-add,inc:req:bre:prop:rep-mult}),
and that any function defined from representable functions by
composition or regular minimization
(\olref[cmp]{prop:rep-composition},
\olref[min]{prop:rep-minimization}) is also representable. Thus every
general recursive function is representable in~$\Th{Q}$.
\end{proof}

\begin{explain}
We have shown that the set of computable functions can be
characterized as the set of functions representable in $\Th{Q}$. In
fact, the proof is more general. From the definition of
representability, it is not hard to see that any theory extending
$\Th{Q}$ (or in which one can interpret $\Th{Q}$) can represent the
computable functions. But, conversely, in any !!{derivation} system in
which the notion of !!{derivation} is computable, every representable
function is computable. So, for example, the set of computable
functions can be characterized as the set of functions representable
in Peano arithmetic, or even Zermelo--Fraenkel set theory. As G\"odel
noted, this is somewhat surprising.  We will see that when it comes to
provability, questions are very sensitive to which theory you
consider; roughly, the stronger the axioms, the more you can prove.
But across a wide range of axiomatic theories, the representable
functions are exactly the computable ones; stronger theories do not
represent more functions as long as they are axiomatizable.
\end{explain}


% ==== OLP file: content/incompleteness/representability-in-q/composition-representable.tex ====
% Part: incompleteness
% Chapter: representability-in-q
% Section: composition-representable



\olfileid{inc}{req}{cmp}
\olsection{Composition is Representable in $\Th{Q}$}

Suppose $h$ is defined by
\[
h(x_0,\dots,x_{l-1}) = f(g_0(x_0,\dots,x_{l-1}), \dots,
g_{k-1}(x_0,\dots,x_{l-1})).
\]
where we have already found !!{formula}s $!A_f, !A_{g_0}, \dots,
!A_{g_{k-1}}$ representing the functions $f$, and $g_0$,
\dots,~$g_{k-1}$, respectively. We have to find !!a{formula}~$!A_h$
representing $h$.

Let's start with a simple case, where all functions are $1$-place,
i.e., consider $h(x) = f(g(x))$. If $!A_f(y, z)$ represents~$f$, and
$!A_g(x, y)$ represents~$g$, we need !!a{formula}~$!A_h(x, z)$ that
represents~$h$.  Note that $h(x) = z$ iff there is a~$y$ such that
both $z = f(y)$ and $y = g(x)$. (If $h(x) = z$, then $g(x)$ is such
a~$y$; if such a $y$ exists, then since $y = g(x)$ and $z = f(y)$, $z
= f(g(x))$.) This suggests that $\lexists[y][(!A_g(x, y) \land !A_f(y,
  z))]$ is a good candidate for~$!A_h(x, z)$. We just have to verify
that $\Th{Q}$ proves the relevant !!{formula}s.

\begin{prop}
\ollabel{prop:rep1}
If $h(n) = m$, then $\Th{Q} \Proves !A_h(\num{n}, \num{m})$.
\end{prop}

\begin{proof}
Suppose $h(n) = m$, i.e., $f(g(n)) = m$.  Let $k = g(n)$. Then
\begin{align*}
  \Th{Q} & \Proves !A_g(\num{n}, \num{k})
  \intertext{since $!A_g$ represents~$g$, and}
  \Th{Q} & \Proves !A_f(\num{k}, \num{m})
  \intertext{since $!A_f$ represents~$f$. Thus,}
  \Th{Q} & \Proves !A_g(\num{n}, \num{k}) \land !A_f(\num{k}, \num{m})
  \intertext{and consequently also}
  \Th{Q} & \Proves \lexists[y][(!A_g(\num{n}, y) \land !A_f(y, \num{m}))],
\end{align*}
i.e., $\Th{Q} \Proves !A_h(\num{n}, \num{m})$.
\end{proof}

\begin{prop}
\ollabel{prop:rep2}
If $h(n) = m$, then $\Th{Q} \Proves \lforall[z][(!A_h(\num{n}, z) \lif
  z = \num{m})]$.
\end{prop}

\begin{proof}
Suppose $h(n) = m$, i.e., $f(g(n)) = m$.  Let $k = g(n)$. Then
\begin{align*}
  \Th{Q} & \Proves \lforall[y][(!A_g(\num{n}, y) \lif \eq[y][\num{k}])]
  \intertext{since $!A_g$ represents~$g$, and}
  \Th{Q} & \Proves \lforall[z][(!A_f(\num{k}, z) \lif \eq[z][\num{m}])]
  \intertext{since $!A_f$ represents~$f$. Using just a little bit of
    logic, we can show that also}
  \Th{Q} & \Proves \lforall[z][(\lexists[y][(!A_g(\num{n}, y) \land
      !A_f(y, z))] \lif \eq[z][\num{m}])].
\end{align*}
i.e., $\Th{Q} \Proves \lforall[y][(!A_h(\num n, y) \lif \eq[y][\num m])]$.
\end{proof}

The same idea works in the more complex case where $f$ and~$g_i$ have
arity greater than~$1$.

\begin{prop}
\ollabel{prop:rep-composition}
If $!A_f(y_0, \dots, y_{k-1}, z)$ represents $f(y_0, \dots, y_{k-1})$
in~$\Th{Q}$, and $!A_{g_i}(x_0, \dots, x_{l-1}, y)$ represents
$g_i(x_0, \dots, x_{l-1})$ in~$\Th{Q}$, then
\begin{multline*}
  \lexists[y_0\dots][\lexists[y_{k-1}][(!A_{g_0}(x_0,\dots,x_{l-1},y_0) \land
      \dots \land {}]]\\
  !A_{g_{k-1}}(x_0,\dots,x_{l-1},y_{k-1}) \land !A_f(y_0,\dots,y_{k-1},z))
\end{multline*}
represents
\[
h(x_0, \dots, x_{l-1}) = f(g_0(x_0, \dots, x_{l-1}), \dots, g_{k-1}(x_0,
\dots, x_{l-1})).
\]
\end{prop}

\begin{proof}
Exercise.
\end{proof}

\begin{prob}
Using the proofs of \olref[inc][req][cmp]{prop:rep2} and
\olref[inc][req][cmp]{prop:rep2} as a guide, carry out the proof of
\olref[inc][req][cmp]{prop:rep-composition} in detail.
\end{prob}


% ==== OLP file: content/incompleteness/representability-in-q/introduction.tex ====
% Part: incompleteness
% Chapter: representability-in-q
% Section: introduction



\olfileid{inc}{req}{int}
\olsection{Introduction}

The incompleteness theorems apply to theories in which basic facts
about computable functions can be expressed and proved.  We will
describe a very minimal such theory called ``$\Th{Q}$'' (or, sometimes,
``Robinson's $Q$,'' after Raphael Robinson). We will say what it means
for a function to be \emph{representable} in $\Th{Q}$, and then we
will prove the following:
\begin{quote}
  A function is representable in $\Th{Q}$ if and only if it is computable.
\end{quote}
For one thing, this provides us with another model of
computability. But we will also use it to show that the set
$\Setabs{!A}{\Th{Q} \Proves !A}$ is not decidable, by reducing the
halting problem to it. By the time we are done, we will have proved
much stronger things than this.

The language of $\Th{Q}$ is the language of
arithmetic; $\Th{Q}$ consists of the following axioms
(to be used in conjunction with the other axioms and rules of
first-order logic with !!{identity}):
\begin{align*}
& \lforall[x][\lforall[y][(\eq[x'][y'] \lif \eq[x][y])]] \tag{$!Q_1$}\\
& \lforall[x][\eq/[\Obj 0][x']] \tag{$!Q_2$}\\
& \lforall[x][(\eq[x][\Obj 0] \lor \lexists[y][\eq[x][y']])] \tag{$!Q_3$}\\
& \lforall[x][\eq[(x + \Obj 0)][x]] \tag{$!Q_4$}\\
& \lforall[x][\lforall[y][\eq[(x + y')][(x + y)']]] \tag{$!Q_5$}\\
& \lforall[x][\eq[(x \times \Obj 0)][\Obj 0]] \tag{$!Q_6$}\\
& \lforall[x][\lforall[y][\eq[(x \times y')][((x \times y) + x)]]] \tag{$!Q_7$}\\
& \lforall[x][\lforall[y][(x < y \liff \lexists[z][\eq[(z' + x)][y]])]] \tag{$!Q_8$}
\end{align*}
For each natural number $n$, define the numeral $\num{n}$ to be the
term $\Obj{0}^{\prime\prime\ldots\prime}$ where there are $n$ tick marks in
all.  So, $\num{0}$ is the !!{constant}~$\Obj{0}$ by itself, $\num{1}$
is $\Obj{0}'$, $\num{2}$ is $\Obj{0}''$, etc.

As a theory of arithmetic, $\Th{Q}$ is \emph{extremely} weak; for
example, you can't even prove very simple facts like
$\lforall[x][\eq/[x][x']]$ or $\lforall[x][\lforall[y][(x + y) = (y +
    x)]]$. But we will see that much of the reason that $\Th{Q}$ is so
interesting is \emph{because} it is so weak. In fact, it is just
barely strong enough for the incompleteness theorem to hold. Another
reason $\Th{Q}$ is interesting is because it has a \emph{finite} set
of axioms.

A stronger theory than $\Th{Q}$ (called \emph{Peano arithmetic} $\Th{PA}$)
is obtained by adding a schema of induction to~$\Th{Q}$:
\[
(!A(\Obj 0) \land \lforall[x][(!A(x) \lif !A(x'))]) \lif \lforall[x][!A(x)]
\]
where $!A(x)$ is any formula. If $!A(x)$ contains free !!{variable}s
other than $x$, we add universal quantifiers to the front to bind all
of them (so that the corresponding instance of the induction schema is
!!a{sentence}). For instance, if $!A(x, y)$ also contains the
!!{variable}~$y$ free, the corresponding instance is
\[
\lforall[y][((!A(\Obj 0) \land \lforall[x][(!A(x) \lif !A(x'))]) \lif
  \lforall[x][!A(x)])]
\]
Using instances of the induction schema, one can prove much more from
the axioms of~$\Th{PA}$ than from those of $\Th{Q}$. In fact, it takes
a good deal of work to find ``natural'' statements about the natural
numbers that can't be proved in Peano arithmetic!{}

\begin{defn}
\ollabel{defn:representable-fn}
  A function $f(x_0,\ldots,x_k)$ from the natural numbers to
  the natural numbers is said to be {\em representable in $\Th{Q}$} if
  there is a formula $!A_f(x_0,\dots,x_k,y)$ such that whenever
  $f(n_0,\dots,n_k) = m$, $\Th{Q}$ proves
\begin{enumerate}
\item\ollabel{defn:rep:a} $!A_f(\num{n_0}, \dots, \num{n_k}, \num{m})$
\item\ollabel{defn:rep:b} $\lforall[y][(!A_f(\num{n_0}, \dots,
\num{n_k}, y) \lif \num{m} = y)]$.
\end{enumerate}
\end{defn}

There are other ways of stating the definition; for example, we could
equivalently require that $\Th{Q}$ proves $\lforall[y][(!A_f(\num{n_0}, \dots,
\num{n_k}, y) \liff \eq[y][\num{m}])]$.

\begin{thm}
\ollabel{thm:representable-iff-comp}
A function is representable in $\Th{Q}$ if and only if it is computable.
\end{thm}

There are two directions to proving the theorem. The left-to-right
direction is fairly straightforward once arithmetization of syntax is
in place. The other direction requires more work.  Here is the basic
idea: we pick ``general recursive'' as a way of making ``computable''
precise, and show that every general recursive function is
representable in~$\Th{Q}$. Recall that a function is general recursive
if it can be defined from $\Zero$, the successor function~$\Succ$, and the
projection functions~$\Proj{n}{i}$, using composition, primitive recursion,
and regular minimization. So one way of showing that every general
recursive function is representable in~$\Th{Q}$ is to show that the
basic functions are representable, and whenever some functions are
representable, then so are the functions defined from them using
composition, primitive recursion, and regular minimization. In other
words, we might show that the basic functions are representable, and
that the representable functions are ``closed under'' composition,
primitive recursion, and regular minimization.  This guarantees that
every general recursive function is representable.

It turns out that the step where we would show that representable
functions are closed under primitive recursion is hard. In order to
avoid this step, we show first that in fact we can do without
primitive recursion. That is, we show that every general recursive
function can be defined from basic functions using composition and
regular minimization alone.  To do this, we show that primitive
recursion can actually be done by a specific regular minimization.
However, for this to work, we have to add some additional basic
functions: addition, multiplication, and the characteristic function
of the identity relation~$\Char{=}$.  Then, we can prove the theorem
by showing that all of \emph{these} basic functions are representable
in~$\Th{Q}$, and the representable functions are closed under
composition and regular minimization.


% ==== OLP file: content/incompleteness/representability-in-q/minimization-representable.tex ====
% Part: incompleteness
% Chapter: representability-in-q
% Section: minimization-representable



\olfileid{inc}{req}{min}
\olsection{Regular Minimization is Representable in $\Th{Q}$}

Let's consider unbounded search. Suppose $g(x, z)$ is regular and
representable in $\Th{Q}$, say by the !!{formula}~$!A_g(x, z, y)$. Let
$f$ be defined by $f(z) = \umin{x}{[g(x, z) = 0]}$. We would like to find
!!a{formula}~$!A_f(z, y)$ representing~$f$.  The value of~$f(z)$ is
that number~$x$ which (a) satisfies $g(x, z) = 0$ and (b) is the least
such, i.e., for any $w < x$, $g(w, z) \neq 0$.  So the following is a
natural choice:
\[
!A_f(z,y) \ident !A_g(y, z, \Obj 0) \land \lforall[w][(w < y \lif \lnot
  !A_g(w, z, \Obj 0))].
\]
In the general case, of course, we would have to replace $z$ with
$z_0$, \dots, $z_k$.

The proof, again, will involve some lemmas about things $\Th{Q}$ is
strong enough to prove.

\begin{lem}
\ollabel{lem:succ} For every !!{constant}~$a$ and every natural
number~$n$,
\[
\Th{Q} \Proves \eq[(a' + \num n)][(a + \num n)'].
\]
\end{lem}

\begin{proof}
The proof is, as usual, by induction on $n$. In the base case, $n =
0$, we need to show that $\Th{Q}$ proves $\eq[(a' + \Obj 0)][(a + \Obj
0)']$. But we have:
\begin{align}
  \Th{Q} & \Proves \eq[(a' + \Obj 0)][a'] \quad \text{by axiom $Q_4$}
  \ollabel{step1}\\
  \Th{Q} & \Proves  \eq[(a + \Obj 0)][a] \quad \text{by axiom $Q_4$}
  \ollabel{step2} \\
  \Th{Q} & \Proves \eq[(a + \Obj 0)'][a'] \quad \text{by \olref{step2}}
  \ollabel{step3} \\
  \Th{Q} & \Proves \eq[(a' + \Obj 0)][(a + \Obj 0)'] \quad
  \text{by \olref{step1} and \olref{step3}}\notag
\end{align}
In the induction step, we can assume that we have shown that $\Th{Q}
\Proves \eq[(a' + \num n)][(a + \num n)']$. Since
$\num{n+1}$ is $\num{n}'$, we need to show that $\Th{Q}$ proves
$\eq[(a' + \num{n}')][(a + \num{n}')']$. We have:
\begin{align}
  \Th{Q} & \Proves \eq[(a' + \num n')][(a' + \num n)'] \quad
  \text{by axiom $!Q_5$} \ollabel{step5}\\
  \Th{Q} & \Proves \eq[(a' + \num n')][(a + \num n')'] \quad
  \text{inductive hypothesis} \ollabel{step6}\\
  \Th{Q} & \Proves \eq[(a' + \num n)'][(a + \num n')'] \quad
  \text{by \olref{step5} and \olref{step6}.} \notag
\end{align}
\end{proof}

It is again worth mentioning that this is weaker than saying that
$\Th{Q}$ proves $\lforall[x][\lforall[y][(x' + y) = (x + y)']]$.
Although this !!{sentence} is true in~$\Struct{N}$, $\Th{Q}$ does not
prove it.

\begin{lem}
\ollabel{lem:less-zero}
$\Th{Q} \Proves \lforall[x][\lnot x < \Obj 0]$.
\end{lem}

\begin{proof}
We give the proof informally (i.e., only giving hints as to
how to construct the formal !!{derivation}).

We have to prove $\lnot a < \Obj 0$ for an arbitrary~$a$. By the
definition of $<$, we need to prove $\lnot
\lexists[y][\eq[(y'+a)][\Obj 0]]$ in $\Th{Q}$. We'll assume
$\lexists[y][\eq[(y'+a)][\Obj 0]]$ and prove a contradiction. Suppose
$\eq[(b'+a)][\Obj 0]$. Using $!Q_3$, we
have that $\eq[a][\Obj 0] \lor \lexists[y][\eq[a][y']]$. We
distinguish cases. 

Case 1: $\eq[a][\Obj 0]$ holds. From $\eq[(b'+a)][\Obj 0]$, we have
$\eq[(b' + \Obj 0)][\Obj 0]$. By axiom~$!Q_4$ of $\Th{Q}$, we have
$\eq[(b' + \Obj 0)][b']$, and hence $\eq[b'][\Obj 0]$. But by
axiom~$!Q_2$ we also have $\eq/[b'][\Obj 0]$, a contradiction. 

Case 2: For some $c$, $\eq[a][c']$. But then we have $\eq[(b' +
c')][\Obj 0]$. By axiom~$!Q_5$, we have $\eq[(b' + c)'][\Obj 0]$, again
contradicting axiom~$Q_2$.
\end{proof}

\begin{lem}
\ollabel{lem:less-nsucc}
For every natural number~$n$, 
  \[
  \Th{Q} \Proves
  \lforall[x][(x < \num {n+1} \lif (\eq[x][\Obj 0] \lor \dots \lor
    \eq[x][\num n]))].
  \]
\end{lem}

\begin{proof}
We use induction on~$n$. Let us consider the base case, when $n = 0$.
In that case, we need to show $a < \num 1 \lif \eq[a][\Obj 0]$, for
arbitrary~$a$. Suppose $a < \num 1$. Then by the defining axiom for
$<$, we have $\lexists[y][\eq[(y'+a)][\Obj 0']]$ (since $\num 1 \ident
\Obj 0'$).

Suppose $b$ has that property, i.e., we have $\eq[(b'+a)][\Obj 0']$.
We need to show $\eq[a][\Obj 0]$. By axiom~$!Q_3$, we have either
$\eq[a][\Obj 0]$ or that there is a $c$ such that $\eq[a][c']$. In the
former case, there is nothing to show. So suppose $\eq[a][c']$. Then
we have $\eq[(b' + c')][\Obj 0']$. By axiom~$!Q_5$ of $\Th{Q}$, we have
$\eq[(b'+c)'][\Obj 0']$. By axiom~$!Q_1$, we have $\eq[(b' + c)][\Obj
0]$. But this means, by axiom~$!Q_8$, that $c < \Obj 0$, contradicting
\olref{lem:less-zero}.

Now for the inductive step. We prove the case for $n+1$, assuming the
case for~$n$. So suppose $a < \num {n+2}$. Again using $!Q_3$ we can
distinguish two cases: $\eq[a][\Obj 0]$ and for some $b$,
$\eq[a][b']$. In the first case, $\eq[a][\Obj 0] \lor \dots \lor
\eq[a][\num{n+1}]$ follows trivially. In the second case, we have $b'
< \num {n+2}$, i.e., $b' < \num{n+1}'$. By axiom~$!Q_8$, for some $c$,
$\eq[(c'+b')][\num{n+1}']$. By axiom $!Q_5$,
$\eq[(c'+b)'][\num{n+1}']$. By axiom~$!Q_1$, $\eq[(c'+b)][\num{n+1}]$,
and so $b < \num{n+1}$ by axiom~$!Q_8$. By inductive hypothesis,
$\eq[b][\Obj 0] \lor \dots \lor \eq[b][\num{n}]$. From this, we get
$\eq[b'][\Obj 0'] \lor \dots \lor \eq[b'][\num{n}']$ by logic, and so
$\eq[a][\num{1}] \lor \dots \lor \eq[a][\num{n+1}]$ since
$\eq[a][b']$.
\end{proof}

\begin{lem}
  \ollabel{lem:trichotomy} For every natural number~$m$,
  \[
  \Th{Q} \Proves
  \lforall[y][((y < \num{m} \lor \num{m} < y) \lor \eq[y][\num{m}])].
  \]
\end{lem}

\begin{proof}
By induction on~$m$. First, consider the case $m=0$. $\Th{Q} \Proves
\lforall[y][(\eq[y][\Obj 0] \lor \lexists[z][\eq[y][z']])]$ by~$!Q_3$.
Let $a$ be arbitrary. Then either $\eq[a][\Obj 0]$ or for some~$b$,
$\eq[a][b']$. In the former case, we also have $(a < \Obj 0 \lor \Obj
0 < a) \lor \eq[a][\Obj 0]$. But if $\eq[a][b']$, then $\eq[(b' +
\Obj 0)][(a + \Obj 0)]$ by the logic of~$\eq$. By $!Q_4$, $\eq[(a +
\Obj 0)][a]$, so we have $\eq[(b' + \Obj 0)][a]$, and hence
$\lexists[z][\eq[(z' + \Obj 0)][a]]$. By the definition of $<$ in
$!Q_8$, $\Obj 0 < a$.  If $\Obj 0 < a$, then also $(\Obj 0 < a \lor a
< \Obj 0) \lor \eq[a][\Obj 0]$. 

Now suppose we have
\begin{align*}
  \Th{Q} & \Proves \lforall[y][((y < \num{m} \lor \num{m} < y) \lor
    \eq[y][\num{m}])]
  \intertext{and we want to show}
  \Th{Q} & \Proves \lforall[y][((y < \num{m+1} \lor \num{m+1} < y) \lor
    \eq[y][\num{m+1}])]
\end{align*}
Let $a$ be arbitrary. By $!Q_3$, either $\eq[a][\Obj 0]$ or for
some~$b$, $\eq[a][b']$. In the first case, we have $\eq[\num{m}' +
a][\num{m+1}]$ by $!Q_4$, and so $a < \num{m+1}$ by $!Q_8$.

Now consider the second case, $\eq[a][b']$. By the induction
hypothesis, $(b < \num{m} \lor \num{m} < b) \lor
    \eq[b][\num{m}]$.

The first disjunct $b < \num{m}$ is equivalent (by $!Q_8$) to
$\lexists[z][\eq[(z' + b)][\num{m}]]$. Suppose $c$ has this property.
If $\eq[(c' + b)][\num{m}]$, then also $\eq[(c' + b)'][\num{m}']$. By
$!Q_5$, $\eq[(c' + b)'][(c' + b')]$. Hence, $\eq[(c' +
b')][\num{m}']$. We get $\lexists[u][\eq[(u' + b')][\num{m+1}]]$ by
existentially generalizing on~$c'$ and keeping in mind that $\num{m}'
\ident \num{m+1}$. Hence, if $b < \num{m}$ then $b' < \num{m+1}$ and
so $a < \num{m+1}$.

Now suppose $\num{m} < b$, i.e., $\lexists[z][\eq[(z' +
\num{m})][b]]$.  Suppose $c$ is such a~$z$, i.e., $\eq[(c' +
\num{m})][b]$. By logic, $\eq[(c' + \num{m})'][b']$. By $!Q_5$,
$\eq[(c' + \num{m}')][b']$. Since $\eq[a][b']$ and $\num{m}' \ident
\num{m+1}$, $\eq[(c' + \num{m+1})][a]$. By $!Q_8$, $\num{m+1} < a$.

Finally, assume $\eq[b][\num{m}]$. Then, by logic,
$\eq[b'][\num{m}']$, and so $\eq[a][\num{m+1}]$.

Hence, from each disjunct of the case for~$m$ and~$b$, we can obtain
the corresponding disjunct for~$m+1$ and~$a$.
\end{proof}
  
\begin{prop}
\ollabel{prop:rep-minimization}
If $!A_g(x, z, y)$ represents $g(x, z)$ in~$\Th{Q}$, then
\[
!A_f(z,y) \ident !A_g(y, z, \Obj 0) \land \lforall[w][(w < y \lif \lnot
  !A_g(w, z, \Obj 0))]
\]
represents $f(z) = \umin{x}{[g(x, z) = 0]}$.
\end{prop}

\begin{proof}
First we show that if $f(n) = m$, then $\Th{Q} \Proves !A_f(\num n, \num m)$,
i.e.,  
\begin{align}
  \Th{Q} & \Proves !A_g(\num{m}, \num{n}, \Obj 0) \land \lforall[w][(w <
    \num{m} \lif \lnot !A_g(w, \num{n}, \Obj 0))]. \notag
  \intertext{Since $!A_g(x, z, y)$ represents $g(x, z)$ and $g(m, n) =
    0$ if $f(n) = m$, we have}
\Th{Q} & \Proves !A_g(\num{m}, \num{n}, \Obj 0). \notag
\intertext{If $f(n) = m$, then for every $k < m$, $g(k, n) \neq 0$. So}
\Th{Q} & \Proves \lnot !A_g(\num{k}, \num{n}, \Obj 0). \notag
\intertext{We get that}
\Th{Q} & \Proves \lforall[w][(w < \num{m} \lif \lnot
  !A_g(w, \num{n}, \Obj 0))]. \ollabel{rep-less}
\end{align}
by \olref{lem:less-zero} in case $m = 0$ and by \olref{lem:less-nsucc} otherwise.

Now let's show that if $f(n) = m$, then $\Th{Q} \Proves
\lforall[y][(!A_f(\num{n}, y) \lif \eq[y][\num{m}])]$.  We again
sketch the argument informally, leaving the formalization to the
reader.

Suppose $!A_f(\num{n}, b)$. From this we get (a) $!A_g(b, \num{n},
\Obj 0)$ and (b) $\lforall[w][(w < b \lif \lnot !A_g(w, \num{n}, \Obj
  0))]$.  By \olref{lem:trichotomy}, $(b < \num{m} \lor \num{m} < b)
\lor \eq[b][\num{m}]$. We'll show that both $b < \num{m}$ and $\num{m}
< b$ leads to a contradiction.

If $\num{m} < b$, then $\lnot !A_g(\num{m}, \num{n}, \Obj 0)$
from~(b). But $m = f(n)$, so $g(m, n) = 0$, and so $\Th{Q} \Proves
!A_g(\num{m}, \num{n}, \Obj 0)$ since $!A_g$ represents~$g$. So we
have a contradiction.

Now suppose $b < \num{m}$. Then since $\Th{Q} \Proves \lforall[w][(w <
  \num{m} \lif \lnot !A_g(w, \num{n}, \Obj 0))]$ by \olref{rep-less}, we
get $\lnot !A_g(b, \num{n}, \Obj 0)$. This again contradicts~(a).
\end{proof}


% ==== OLP file: content/incompleteness/representability-in-q/prim-rec.tex ====
% Part: incompleteness
% Chapter: representability-in-q
% Section: prim-rec



\olfileid{inc}{req}{pri}
\olsection{Simulating Primitive Recursion}

Now we can show that definition by primitive recursion can be
``simulated'' by regular minimization using the beta function. Suppose
we have $f(\vec x)$ and $g(\vec x, y, z)$. Then the function~$h(x,\vec
z)$ defined from $f$ and~$g$ by primitive recursion is
\begin{align*}
h(\vec x, 0) & =  f(\vec x) \\
h(\vec x, y+1) & =  g(\vec x, y, h(\vec x, y)).
\end{align*}
We need to show that $h$ can be defined from $f$ and~$g$ using just
composition and regular minimization, using the basic functions and
functions defined from them using composition and regular minimization
(such as~$\beta$).

\begin{lem}
\ollabel{lem:prim-rec}
If $h$ can be defined from $f$ and $g$ using primitive recursion, it
can be defined from $f$, $g$, the functions $\Zero$, $\Succ$,
$\Proj{n}{i}$, $\Add$, $\Mult$, $\Char{=}$, using composition and
regular minimization.
\end{lem}

\begin{proof}
First, define an auxiliary function $\hat h(\vec x, y)$ which returns
the least number~$d$ such that $d$ codes a sequence which satisfies
\begin{enumerate}
\item $(d)_0 = f(\vec x)$, and
\item for each $i < y$, $(d)_{i+1} = g(\vec x, i, (d)_i)$,
\end{enumerate}
where now $(d)_i$ is short for $\beta(d,i)$. In other words, $\hat h$
returns the sequence $\tuple{h(\vec x, 0), h(\vec x, 1), \dots, h(\vec
x, y)}$. We can write $\hat h$ as
\[
\hat h(\vec x, y) = \umin{d}{(\beta(d,0) = f(\vec x) \land \bforall{i <
  y}{\beta(d,i+1) = g(\vec x, i,\beta(d,i)})}.
\]
Note: no primitive recursion is needed here, just minimization. The
function we minimize is regular because of the beta function lemma
\olref[bet]{lem:beta}.

But now we have
\[
h(\vec x, y) = \beta(\hat h(\vec x, y), y),
\]
so $h$ can be defined from the basic functions using just composition
and regular minimization.
\end{proof}


% ==== OLP file: content/incompleteness/representability-in-q/representability-in-q.tex ====
% Part: incompleteness
% Chapter: representability-in-q



\olchapter{inc}{req}{Representability in $\Th{Q}$}

\olimport{introduction}

\olimport{representable-comp}

\olimport{beta-function}

\olimport{prim-rec}

\olimport{basic-representable}

\olimport{composition-representable}

\olimport{minimization-representable}

\olimport{comp-representable}

\olimport{representing-relations}

\olimport{undecidability}

\olimport{sigma1-completeness}

\OLEndChapterHook


% ==== OLP file: content/incompleteness/representability-in-q/representable-comp.tex ====
% Part: incompleteness
% Chapter: representability-in-q
% Section: representable-comp



\olfileid{inc}{req}{rpc}
\olsection{Functions Representable in $\Th{Q}$ are Computable}

We'll prove that every function that is representable in~$\Th{Q}$ is
computable. We first have to establish a lemma about functions
representable in~$\Th{Q}$.

\begin{lem}\ollabel{lem:rep-q}
  If $f(x_0,
  \dots, x_k)$ is representable in~$\Th{Q}$, there is
  !!a{formula}~$!A(x_0, \dots, x_k, y)$ such that
  \[
  \Th{Q} \Proves !A_f(\num{n_0}, \dots, \num{n_k}, \num{m})
  \quad\text{iff}\quad m = f(n_0, \dots, n_k).
  \]
\end{lem}

\begin{proof}
  The ``if'' part is
\olref[int]{defn:representable-fn}\olref[int]{defn:rep:a}. The ``only
if'' part is seen as follows: Suppose $\Th{Q} \Proves !A_f(\num{n_0},
\dots, \num{n_k}, \num{m})$ but $m \neq f(n_0, \dots, n_k)$. Let $l =
f(n_0, \dots, n_k)$. By
\olref[int]{defn:representable-fn}\olref[int]{defn:rep:a}, $\Th{Q}
\Proves !A_f(\num{n_0}, \dots, \num{n_k}, \num{l})$. By
\olref[int]{defn:representable-fn}\olref[int]{defn:rep:b},
$\lforall[y][(!A_f(\num{n_0}, \dots, \num{n_k}, y) \lif \num{l} =
y)]$. Using logic and the assumption that $\Th{Q} \Proves
!A_f(\num{n_0}, \dots, \num{n_k}, \num{m})$, we get that $\Th{Q}
\Proves \eq[\num{l}][\num{m}]$. On the other hand, by
\olref[bre]{lem:q-proves-neq}, $\Th{Q} \Proves
\eq/[\num{l}][\num{m}]$. So $\Th{Q}$ is inconsistent. But that is
impossible, since $\Th{Q}$ is satisfied by the standard model (see
\olref[int][def]{def:standard-model}), $\Sat{N}{\Th{Q}}$, and
satisfiable theories are always consistent by the Soundness Theorem
(\tagrefs{prfAX/{fol:axd:sou:cor:consistency-soundness},
prfSC/{fol:seq:sou:cor:consistency-soundness},
prfND/{fol:ntd:sou:cor:consistency-soundness},
prfTab/{fol:tab:sou:cor:consistency-soundness}}).
\end{proof}

\begin{lem}
Every function that is representable in $\Th{Q}$ is computable.
\end{lem}

\begin{proof}
Let's first give the intuitive idea for why this is true. To
compute~$f$, we do the following.  List all the possible
!!{derivation}s~$\delta$ in the language of arithmetic. This is
possible to do mechanically. For each one, check if it is
!!a{derivation} of !!a{formula} of the form~$!A_f(\num{n_0}, \dots,
\num{n_k}, \num{m})$ (the !!{formula} representing $f$ in~$\Th{Q}$
from \olref{lem:rep-q}). If it is, $m = f(n_0, \dots, n_k)$ by
\olref{lem:rep-q}, and we've found the value of~$f$. The search
terminates because $\Th{Q} \Proves !A_f(\num{n_0}, \dots, \num{n_k},
\num{f(n_0, \dots, n_k)})$, so eventually we find a $\delta$ of the
right sort.

This is not quite precise because our procedure operates on
!!{derivation}s and !!{formula}s instead of just on numbers, and we
haven't explained exactly why ``listing all possible !!{derivation}s''
is mechanically possible.  But as we've seen, it is possible to code
terms, !!{formula}s, and !!{derivation}s by G\"odel numbers. We've
also introduced a precise model of computation, the general recursive
functions. And we've seen that the relation $\Prf[\Th{Q}](d,y)$, which
holds iff $d$ is the G\"odel number of !!a{derivation} of the !!{formula}
with G\"odel number~$y$ from the axioms of~$\Th{Q}$, is (primitive)
recursive. Other primitive recursive functions we'll need are
$\fn{num}$ (\olref[art][trm]{prop:num-primrec}) and $\fn{Subst}$
(\olref[art][sub]{prop:subst-primrec}).  From these, it is possible to
define~$f$ by minimization; thus, $f$ is recursive.

First, define
\begin{multline*}
  A(n_0, \dots, n_k, m) = \\
  \fn{Subst}(\fn{Subst}(\dots\fn{Subst}(\Gn{!A_f}, \fn{num}(n_0), \Gn{x_0}),\\ \dots),
  \fn{num}(n_k),  \Gn{x_k}), \fn{num}(m), \Gn{y})
\end{multline*}
This looks complicated, but it's just the function $A(n_0, \dots, n_k,
m) = \Gn{!A_f(\num{n_0}, \dots, \num{n_k}, \num{m})}$.

Now, consider the relation~$R(n_0, \dots, n_k, s)$ which holds if
$(s)_0$ is the G\"odel number of !!a{derivation} from~$\Th{Q}$ of
$!A_f(\num{n_0}, \dots, \num{n_k}, \num{(s)_1})$:
\[
R(n_0, \dots, n_k, s) \quad\text{iff}\quad \Prf[\Th{Q}]((s)_0, A(n_0,
\dots, n_k, (s)_1))
\]
If we can find an~$s$ such that $R(n_0, \dots, n_k, s)$ holds, we have
found a pair of numbers---$(s)_0$ and~$(s)_1$---such that $(s)_0$ is
the G\"odel number of !!a{derivation} of~$A_f(\num{n_0}, \dots,
\num{n_k}, (s)_1)$. So looking for~$s$ is like looking for the pair
$d$ and $m$ in the informal proof. And a computable function that
``looks for'' such an $s$ can be defined by regular minimization.
Note that $R$ is regular: for every $n_0$, \dots, $n_k$, there is
!!a{derivation}~$\delta$ of $\Th{Q} \Proves !A_f(\num{n_0}, \dots,
\num{n_k}, \num{f(n_0, \dots, n_k)})$, so $R(n_0, \dots, n_k, s)$
holds for $s = \tuple{\Gn{\delta}, f(n_0, \dots, n_k)}$.  So, we can
write $f$ as
\[
f(n_0,\dots,n_{k}) = (\umin{s}{R(n_0, \dots, n_k, s)})_1.
\]
\end{proof}


% ==== OLP file: content/incompleteness/representability-in-q/representing-relations.tex ====
% Part: incompleteness
% Chapter: representability-in-q
% Section: representing-relations



\olfileid{inc}{req}{rel}

\olsection{Representing Relations}

Let us say what it means for a \emph{relation} to be representable.

\begin{defn}
\ollabel{defn:representing-relations} A relation $R(x_0,\dots,x_k)$ on
the natural numbers is {\em representable in $\Th{Q}$} if there is a
formula $!A_R(x_0,\dots,x_k)$ such that whenever $R(n_0,\dots,n_k)$ is
true, $\Th{Q}$ proves $!A_R(\num{n_0},\dots,\num{n_k})$, and whenever
$R(n_0,\dots,n_k)$ is false, $\Th{Q}$ proves $\lnot !A_R(\num{n_0},
\dots, \num{n_k})$.
\end{defn}

\begin{thm}
\ollabel{thm:representing-rels} A relation is representable in
$\Th{Q}$ if and only if it is computable.
\end{thm}

\begin{proof}
For the forwards direction, suppose $R(x_0,\dots,x_k)$ is
represented by the formula $!A_R(x_0,\dots,x_k)$. Here is an
algorithm for computing $R$: on input $n_0$, \dots,~$n_k$, simultaneously
search for a proof of $!A_R(\num{n_0}, \dots, \num{n_k})$ and a proof of
$\lnot !A_R(\num{n_0}, \dots, \num{n_k})$. By our hypothesis, the search
is bound to find one or the other; if it is the first, report ``yes,''
and otherwise, report ``no.''

In the other direction, suppose $R(x_0, \dots, x_k)$ is computable. By
definition, this means that the function $\Char{R}(x_0, \dots, x_k)$
is computable. By \olref[int]{thm:representable-iff-comp}, $\Char{R}$
is represented by a formula, say $!A_{\Char{R}}(x_0, \dots, x_k,
y)$. Let $!A_R(x_0, \dots, x_k)$ be the formula $!A_{\Char{R}}(x_0,
\dots, x_k, \num{1})$. Then for any $n_0$, \dots,~$n_k$, if $R(n_0,
\dots, n_k)$ is true, then $\Char{R}(n_0, \dots, n_k) = 1$, in which
case $\Th{Q}$ proves $!A_{\Char{R}}(\num{n_0}, \dots, \num{n_k},
\num{1})$, and so $\Th{Q}$ proves $!A_R(\num{n_0}, \dots,
\num{n_k})$. On the other hand, if $R(n_0, \dots, n_k)$ is false, then
$\Char{R}(n_0, \dots, n_k) = 0$. This means that $\Th{Q}$ proves
\[
\lforall[y][(!A_{\Char{R}}(\num{n_0}, \dots, \num{n_k}, y) \lif y =
  \num{0})].
\]
Since $\Th{Q}$ proves $\eq/[\num{0}][\num{1}]$, $\Th{Q}$ proves
$\lnot !A_{\Char{R}}(\num{n_0}, \dots, \num{n_k}, \num{1})$, and so it
proves $\lnot !A_R(\num{n_0}, \dots, \num{n_k})$.
\end{proof}

\begin{prob}
Show that if $R$ is representable in~$\Th{Q}$, so is~$\Char{R}$.
\end{prob}


% ==== OLP file: content/incompleteness/representability-in-q/sigma1-completeness.tex ====
% Part: incompleteness
% Chapter: representability-in-q
% Section: sigma1-completeness



\olfileid{inc}{inp}{s1c}
\olsection{\texorpdfstring{$\Sigma_1$}{Sigma-1} completeness}

Despite the incompleteness of $\Th{Q}$ and its consistent, axiomatizable
extensions, we have seen that $\Th{Q}$ does prove many basic facts about
numerals. In fact, this can be extended quite considerably. To understand
the scope of what can be proved in~$\Th{Q}$, we introduce the notions of
$\Delta_0$, $\Sigma_1$, and $\Pi_1$ !!{formula}s. Roughly speaking, a
$\Sigma_1$ !!{formula} is one of the form $\lexists[x][!B(x)]$, where $!B$
is constructed using only propositional connectives and bounded
quantifiers. We shall show that if $!A$ is a $\Sigma_1$ !!{sentence}
which is true in $\Struct{N}$, then $\Th{Q} \Proves !A$
(\olref{thm:sigma1-completeness}).

\begin{defn}
\ollabel{defn:bd-quant}
A \emph{bounded existential !!{formula}} is one of the form
$\lexists[x][(x < t \land !A(x))]$ where $t$ is any term, which we
conventionally write as $\bexists{x < t}{!A(x)}$.
%
A \emph{bounded universal !!{formula}} is one of the form
$\lforall[x][(x < t \lif !A(x))]$ where $t$ is any term, which we
conventionally write as $\bforall{x < t}{!A(x)}$.
\end{defn}

\begin{defn}
\ollabel{defn:delta0-sigma1-pi1-frm}
A !!{formula} $!B$ is $\Delta_0$ if it is built up from atomic
!!{formula}s using only propositional connectives and bounded
quantification.
%
A !!{formula} $!A$ is $\Sigma_1$ if $!A \ident \lexists[x][!B(x)]$
where $!B$ is $\Delta_0$.
%
A !!{formula} $!A$ is $\Pi_1$ if $!A \ident \lforall[x][!B(x)]$
where $!B$ is $\Delta_0$.
\end{defn}

\begin{lem}
\ollabel{lem:q-proves-clterm-id} Suppose $t$ is a closed term such that
$\Value{t}{N} = n$. Then $\Th{Q} \Proves \eq[t][\num n]$.
\end{lem}

\begin{proof}
We prove this by induction on the complexity of~$t$. For the base case,
$\Value{\Obj 0}{N} = 0$, and $\Th{Q} \Proves \eq[\Obj 0][\num 0]$
since $\num 0 \ident \Obj 0$.
%
For the inductive case, let $t_1$ and $t_2$ be terms such that
$\Value{t_1}{N} = n_1$, $\Value{t_2}{N} = n_2$,
$\Th{Q} \Proves \eq[t_1][\num n_1]$, and
$\Th{Q} \Proves \eq[t_2][\num n_2]$.

Then $\Value{(t_1')}{N} = n_1 + 1$, and we have that $\Th{Q} \Proves
\eq[t_1'][{\num n_1}']$ by the first-order rules for identity applied
to the induction hypothesis and the !!{formula}
$\eq[\num{n_1}'][\num{n_1}']$,
so we have $\Th{Q} \Proves \eq[t_1'][\num{n_1 + 1}]$
by the definition of numerals.

For sums we have
\[
      \Value{(t_1 + t_2)}{N}
    = \Value{t_1}{N} + \Value{t_2}{N}
    = n_1 + n_2.
\]
By the induction hypothesis and the rules for identity,
$\Th{Q} \Proves \eq[t_1 + t_2][\num{n_1} + t_2]$, and then
$\Th{Q} \Proves \eq[t_1 + t_2][\num{n_1} + \num{n_2}]$
by a second application of the rules for identity.
By \olref[inc][req][bre]{lem:q-proves-add},
$\Th{Q} \Proves \eq[\num{n_1} + \num{n_2}][\num{n_1 + n_2}]$,
so $\Th{Q} \Proves \eq[t_1 + t_2][\num{n_1 + n_2}]$.

Similar reasoning also works for~$\times$, using
\olref[inc][req][bre]{lem:q-proves-mult}.
%
Since this exhausts the closed terms of arithmetic, we have that
$\Th{Q} \Proves \eq[t][\num n]$ for all closed terms~$t$ such that
$\Value{t}{N} = n$.
\end{proof}

\begin{prob}
Prove in detail the part of \olref{lem:q-proves-clterm-id}
involving~$\times$.
\end{prob}

\begin{lem}
\ollabel{lem:atomic-completeness}
Suppose $t_1$ and $t_2$ are closed terms. Then
\begin{enumerate}
\item If $\Value{t_1}{N} = \Value{t_2}{N}$,
    then $\Th{Q} \Proves \eq[t_1][t_2]$.
\item If $\Value{t_1}{N} \neq \Value{t_2}{N}$,
    then $\Th{Q} \Proves \eq/[t_1][t_2]$.
\item If $\Value{t_1}{N} < \Value{t_2}{N}$,
    then $\Th{Q} \Proves t_1 < t_2$.
\item If $\Value{t_2}{N} \leq \Value{t_1}{N}$,
    then $\Th{Q} \Proves \lnot(t_1 < t_2)$.
\end{enumerate}
\end{lem}

\begin{proof}
Given terms $t_1$ and $t_2$, we fix $n = \Value{t_1}{N}$ and
$m = \Value{t_2}{N}$.

Suppose $!A \ident t_1 = t_2$. By \olref{lem:q-proves-clterm-id},
$\Th{Q} \Proves \eq[t_1][\num n]$ and $\Th{Q} \Proves \eq[t_2][\num n]$.
If $n = m$, then $\Th{Q} \Proves \eq[\num n][\num m]$ and hence
$\Th{Q} \Proves \eq[t_1][t_2]$ by the transitivity of identity.
If $n \neq m$ then $\Th{Q} \Proves \eq/[\num n][\num m]$,
and by the transitivity of identity again,
$\Th{Q} \Proves \eq/[t_1][t_2]$.

Now let $!A \ident t_1 < t_2$. For both cases, we rely on axiom~$!Q_8$,
which states that $x < y \liff \lexists[z][\eq[z' + x][y]]$
for all $x,y$.

Suppose $\Sat{N}{t_1 < t_2}$. Then there exists some $k \in \Nat$
such that $n + k + 1 = m$. By \olref{lem:q-proves-clterm-id},
$\Th{Q} \Proves \eq[t_1][\num n]$ and $\Th{Q} \Proves \eq[t_2][\num m]$,
and by the first part of this lemma,
$\Th{Q} \Proves \eq[\num n + {\num k}'][\num m]$.
By the transitivity of identity it follows that
$\Th{Q} \Proves \eq[{\num k}' + t_1][t_2]$,
so $\Th{Q} \Proves \lexists[z][\eq[z' + t_1][t_2]]$.
By the right-to-left direction of~$!Q_8$, $\Th{Q} \Proves t_1 < t_2$.

Suppose instead that $\Sat/{N}{t_1 < t_2}$, i.e., $m \leq n$.
%
We work in~$\Th{Q}$ and assume that $t_1 < t_2$. By the left-to-right
direction of~$!Q_8$, there is some~$z$ such that $\eq[z' + t_1][t_2]$.
Since $\Th{Q} \Proves \eq[t_1][\num n]$ and
$\Th{Q} \Proves \eq[t_2][\num m]$, $\eq[z' + \num n][\num m]$.
%
By an external induction on~$m$ using~$!Q_5$,
$\eq[z' + \num{n - m}][\Obj 0]$.
If $m = n$ then $\eq/[z'][\Obj 0]$, giving a contradiction via~$!Q_3$.
If $m < n$ then $\eq[(z' + \num{n - m - 1})'][\Obj 0]$ by~$!Q_5$ again,
giving a contradiction via~$!Q_3$.
So $\Th{Q} \Proves \lnot(t_1 < t_2)$.
\end{proof}

\begin{lem}
\ollabel{lem:bounded-quant-equiv}
Suppose $!A$ is !!a{formula}, $t$ a closed term, and $k=\Value{t}{N}$. Then
\begin{enumerate}
\item $\Th{Q} \Proves \bforall{x<t}{!A(x)}$ iff $\Th{Q} \Proves
    !A(\num 0) \land \dots \land !A(\num{k-1})$.
\item $\Th{Q} \Proves \bexists{x<t}{!A(x)}$ iff $\Th{Q} \Proves
    !A(\num 0) \lor \dots \lor !A(\num{k-1})$.
\end{enumerate}
\end{lem}

\begin{proof}
    We prove the case for the bounded universal quantifier.
    If $\Value{t}{N} = 0$ then the left-hand side of the
    equivalence is provable in~$\Th{Q}$, because there is no
    $x<\num 0$ by \olref[inc][req][min]{lem:less-zero}.
    Similarly, we can take an empty disjunction to be simply
    $\ltrue$, which is also provable in~$\Th{Q}$.
    %
    We therefore suppose that $\Value{t}{N} = k+1$ for some
    natural number~$k$. By \olref{lem:q-proves-clterm-id} we
    can assume that we are working with !!a{formula} of the
    form $\bforall{x<\num{k+1}}{!A(x)}$.
    
    Suppose that $\Th{Q} \Proves \bforall{x<\num{k+1}}{!A(x)}$,
    and let $n \leq k$. Since $\Th{Q} \Proves \num n < \num{k+1}$
    by \olref{lem:atomic-completeness}, it follows by logic that
    $\Th{Q} \Proves !A(\num n)$. Applying this fact $k+1$ times
    for each $n \leq k$, we get that $\Th{Q} \Proves !A(\num 0)
    \land \dots \land !A(\num k)$ as desired.
    
    For the other direction, suppose that $\Th{Q} \Proves
    !A(\num 0) \land \dots \land !A(\num k)$. Working in
    $\Th{Q}$, suppose that $x < \num{k+1}$.
    By \olref[inc][req][min]{lem:less-nsucc} we have that
    $x = \num 0 \lor \dots \lor x = \num k$, so by logic it
    follows that~$!A(x)$, and hence the universal claim
    $\bforall{x<\num{k+1}}{!A(x)}$ follows.
    
    The proof of the equivalence for bounded existentially
    quantified !!{formula}s is similar.
\end{proof}

\begin{prob}
Give a detailed proof of the existential case in
\olref{lem:bounded-quant-equiv}.
\end{prob}

\begin{lem}
\ollabel{lem:delta0-completeness}
If $!A$ is a $\Delta_0$ !!{sentence} which is true in
$\Struct{N}$, then $\Th{Q} \Proves !A$.
\end{lem}

\begin{proof}
We prove this by induction on !!{formula} complexity.
%
The base case is given by \olref{lem:atomic-completeness},
so we move to the induction step. For simplicity we split
the case of negation into subcases depending on the
structure of the !!{formula} to which the negation is
applied.

\begin{enumerate}
\item Suppose $(!A \land !B)$ is true in $\Struct{N}$,
so $!A$ and $!B$ are true in~$\Struct{N}$.
By the induction hypothesis, $\Th{Q} \Proves !A$ and
$\Th{Q} \Proves !B$,
so $\Th{Q} \Proves (!A \land !B)$ by logic.
%
\item Suppose $\lnot (!A \land !B)$ is true in $\Struct{N}$,
so either $\lnot !A$ or $\lnot !B$ is true in $\Struct{N}$.
Without loss of generality, suppose the former. By the
induction hypothesis $\Th{Q} \Proves \lnot !A$, and hence
$\Th{Q} \Proves \lnot (!A \land !B)$ by logic.
%
\item Suppose $(!A \lor !B)$ is true in $\Struct{N}$, so
either $!A$ is true in $\Struct{N}$ or $!B$ is true in
$\Struct{N}$. Without loss of generality, suppose the former
holds. By the induction hypothesis $\Th{Q} \Proves !A$, and
hence $\Th{Q} \Proves (!A \lor !B)$ by logic.
%
\item Suppose $\lnot(!A \lor !B)$ is true in $\Struct{N}$,
so $\lnot !A$ and $\lnot !B$ are true in $\Struct{N}$.
Then $\Th{Q} \Proves \lnot !A$ and $\Th{Q} \Proves \lnot !B$
by the induction hypothesis. Consequently,
$\Th{Q} \Proves \lnot(!A \lor !B)$ by logic.
%
\item Suppose that $\bforall{x<t}{!A(x)}$ is true 
in~$\Struct{N}$, where $t$ is a closed term and $k=\Value{t}{N}$. By the induction
hypothesis and logic, if $!A(\num n)$ is true in~$\Struct{N}$
for all $n < \Value{t}{N}$ then $\Th{Q} \Proves
!A(\num 0) \land \dots \land !A(\num{k-1})$.
By \olref{lem:bounded-quant-equiv} it follows that
$\Th{Q} \Proves \bforall{x<t}{!A(x)}$.
%
\item The case for the bounded existential quantifier, where
we have !!a{sentence} of the form $\bexists{x < t}{!A(x)}$,
is similar to that for the bounded universal quantifier.
%
\item Suppose that $\lnot \bforall{x<t}{!A(x)}$ is true 
in~$\Struct{N}$, where $t$ is a closed term. This !!{sentence}
is equivalent to the !!{sentence} $\bexists{x<t}{\lnot !A(x)}$,
with the equivalence derivable in~$\Th{Q}$, so we may apply
the reasoning for bounded existential quantifiers.
%
\item Similarly, suppose that $\lnot \bexists{x<t}!A(x)$ is
true in $\Struct{N}$, where $t$ is a closed term. This
!!{sentence} is equivalent in $\Th{Q}$ to
$\bforall{x<t}{\lnot!A(x)}$, and so we may apply the reasoning
for bounded universal quantifiers.
%
\item Finally, suppose $\lnot !A$ is true in $\Struct{N}$.
The only cases remaining are when $!A$ is atomic and when
$\lnot !A \ident \lnot\lnot !B$ for some $\Delta_0$
!!{sentence} $!B$. If $!A$ is atomic then by
\olref{lem:atomic-completeness}, $\Th{Q} \Proves \lnot !A$.
If $\lnot !A \ident \lnot\lnot !B$, then by logic it is
provably equivalent in~$\Th{Q}$ to~$!B$, which is true 
in~$\Struct{N}$ since $\lnot !A$ is true in~$\Struct{N}$.
By the induction hypothesis we therefore have that
$\Th{Q} \Proves \lnot !A$.
\end{enumerate}
\end{proof}

\begin{prob}
Give a detailed proof of the existential case in
\olref{lem:delta0-completeness}.
\end{prob}

\begin{thm}
\ollabel{thm:sigma1-completeness}
If $!A$ is a $\Sigma_1$ !!{sentence} which is true
in~$\Struct{N}$, then $\Th{Q} \Proves !A$.
\end{thm}

\begin{proof}
If $\lexists{x}!A(x)$ is a $\Sigma_1$ !!{sentence} which
is true in~$\Struct{N}$, then there exists a natural
number~$n$ and a variable assignment~$s$ such that $s(x) = n$ and
$\Sat{N}{!A(x)}[s]$. By standard facts about
the satisfaction relation it follows that
$\Sat{N}{!A(\num n)}$. But $!A(\num n)$ is a
$\Delta_0$ !!{formula}, so by \olref{lem:delta0-completeness}
we have that $\Th{Q} \Proves !A(\num n)$, and hence by
logic we also have that $\Th{Q} \Proves \lexists[x][!A(x)]$.
\end{proof}


% ==== OLP file: content/incompleteness/representability-in-q/undecidability.tex ====
% Part: incompleteness
% Chapter: representability-in-q
% Section: undecidability



\olfileid{inc}{req}{und}
\olsection{Undecidability}

We call a theory $\Th{T}$ \emph{undecidable} if there is no
computational procedure which, after finitely many steps and
unfailingly, provides a correct answer to the question ``does $\Th{T}$
prove~$!A$?'' for any sentence~$!A$ in the language of~$\Th{T}$.  So
$\Th{Q}$ would be decidable iff there were a computational procedure
which decides, given a sentence~$!A$ in the language of arithmetic,
whether $\Th{Q} \Proves !A$ or not.  We can make this more precise by
asking: Is the relation~$\Prov[\Th{Q}](y)$, which holds of~$y$
iff $y$ is the G\"odel number of a sentence provable in~$\Th{Q}$,
recursive?  The answer is: no.

\begin{thm}
$\Th{Q}$ is undecidable, i.e., the relation
\[
\Prov[\Th{Q}](y) \defiff \fn{Sent}(y) \land
\lexists[x][\Prf[\Th{Q}](x, y)]
\]
is not recursive.
\end{thm}

\begin{proof}
Suppose it were.  Then we could solve the halting problem as follows:
Given $e$ and $n$, we know that $\cfind{e}(n) \fdefined$ iff there is
an~$s$ such that $T(e, n, s)$, where $T$ is Kleene's predicate from
\olref[cmp][rec][nft]{thm:kleene-nf}.  Since $T$ is primitive recursive
it is representable in~$\Th{Q}$ by a formula $!B_T$, that is, $\Th{Q}
\Proves !B_T(\num{e}, \num{n}, \num{s})$ iff $T(e, n, s)$.  If $\Th{Q}
\Proves !B_T(\num{e}, \num{n}, \num{s})$ then also $ \Th{Q} \Proves
\lexists[y][!B_T(\num{e}, \num{n}, y)]$.  If no such $s$ exists, then
$\Th{Q} \Proves \lnot !B_T(\num{e}, \num{n}, \num{s})$ for
every~$s$.  But $\Th{Q}$ is $\omega$-consistent, i.e., if $\Th{Q}
\Proves \lnot !A(\num{n})$ for every~$n \in \Nat$, then $\Th{Q}
\Proves/ \lexists[y][!A(y)]$.  We know this because the axioms of
$\Th{Q}$ are true in the standard model~$\Struct{N}$.  So, $\Th{Q}
\Proves/ \lexists[y][!B_T(\num{e}, \num{n}, y)]$.  In other words,
$\Th{Q} \Proves \lexists[y][!B_T(\num{e}, \num{n}, y)]$ iff there is
an $s$ such that $T(e, n, s)$, i.e., iff $\cfind{e}(n) \fdefined$.
From $e$ and~$n$ we can compute $\Gn{\lexists[y][!B_T(\num{e},
    \num{n}, y)]}$, let $g(e, n)$ be the primitive recursive function
which does that.  So
\[
h(e, n) =
\begin{cases}
1 & \text{if $\Prov[\Th{Q}](g(e, n))$}\\
0 & \text{otherwise}.
\end{cases}
\]
This would show that $h$ is recursive if $\Prov[\Th{Q}]$ is. But~$h$
is not recursive, by \olref[cmp][rec][hlt]{thm:halting-problem}, so
$\Prov[\Th{Q}]$ cannot be either.
\end{proof}

\begin{cor}
First-order logic is undecidable.
\end{cor}

\begin{proof}
If first-order logic were decidable, provability in~$\Th{Q}$ would be
as well, since $\Th{Q} \Proves !A$ iff $\Proves !T \lif !A$, where
$!T$ is the conjunction of the axioms of~$\Th{Q}$.
\end{proof}

\end{document}

\end{document}

\end{document}
